% =====================================================================
%  Segmented Dealerless Function Secret Sharing for
%  Mixed-Width Private Inference
%
%  Style: IEEEtran (journal), mirroring the FSS-DT template.
%  Build: pdflatex main && bibtex main && pdflatex main && pdflatex main
%         (or: latexmk -pdf main.tex)
%  Requires: IEEEtran.cls, amsmath, amssymb, algorithm, algpseudocode.
% =====================================================================
\documentclass[journal]{IEEEtran}

\usepackage{amsmath,amssymb,amsfonts}
\usepackage{amsthm}
\usepackage{mathtools}
\usepackage{algorithm}
\usepackage{stmaryrd}
\usepackage{algpseudocode}
\usepackage{booktabs}
\usepackage{array}
\usepackage{algorithmicx}
\usepackage{multirow}
\usepackage{graphicx}
% \graphicspath{{./}{paper/}} % resolve figures whether compiled from paper/ or the workspace root
\usepackage{tikz}
\usetikzlibrary{arrows.meta,positioning,fit,backgrounds}
\usepackage{xcolor}
\usepackage{pifont} % \ding{55} (cross mark) in the comparison table
\usepackage[hidelinks]{hyperref}
\usepackage{url}
\usepackage{cite}

% ---------- theorem environments ----------
\newtheorem{definition}{Definition}
\newtheorem{lemma}{Lemma}
\newtheorem{theorem}{Theorem}
\newtheorem{corollary}{Corollary}
\newtheorem{proposition}{Proposition}
\theoremstyle{definition}
\newtheorem{example}{Example}
\theoremstyle{plain}

% ---------- notation macros ----------
\newcommand{\Z}{\mathbb{Z}}
\newcommand{\F}{\mathbb{F}}
\newcommand{\G}{\mathbb{G}}
\newcommand{\R}{\mathcal{R}}                              % ring-generic carrier ring
\newcommand{\ashare}[1]{\langle #1 \rangle}             % arithmetic share
\newcommand{\bshare}[1]{\llbracket #1 \rrbracket}        % boolean share
\newcommand{\rshare}[1]{\langle\!\langle #1 \rangle\!\rangle} % ring-generic share
\newcommand{\Gen}{\mathsf{Gen}}
\newcommand{\Eval}{\mathsf{Eval}}
\newcommand{\DPF}{\mathsf{DPF}}
\newcommand{\DCF}{\mathsf{DCF}}
\newcommand{\Convert}{\mathsf{Convert}}
\newcommand{\Fcot}{\mathcal{F}_{\mathsf{COT}}}
\newcommand{\Frand}{\mathcal{F}_{\mathsf{Rand}}}
\newcommand{\Fmul}{\mathcal{F}_{\mathsf{mul}}}
\newcommand{\Ftr}{\mathcal{F}_{\mathsf{trunc}}}
\newcommand{\hi}{\mathsf{hi}}
\newcommand{\lo}{\mathsf{lo}}
\newcommand{\msb}{\mathsf{msb}}
\newcommand{\lsb}{\mathsf{lsb}}
\newcommand{\DReLU}{\mathsf{DReLU}}
\newcommand{\ReLU}{\mathsf{ReLU}}
\newcommand{\BitBA}{\mathsf{BitBA}}
\newcommand{\ind}[1]{\mathbf{1}\{#1\}}
\newcommand{\xhat}{\hat{x}}
\newcommand{\PREP}{\Pi_{\mathsf{PREP}}}
\newcommand{\CCMP}{\Pi_{\mathsf{CCMP}}}
\newcommand{\NextLevel}{\textsf{NextLevel}}
\newcommand{\bigO}{\mathcal{O}}
\newcommand{\seg}[1]{^{\scriptscriptstyle(#1)}}   % segment index superscript
\newcommand{\nseg}{c}                              % number of segments
\newcommand{\wseg}{w}                              % segment width
\newcommand{\Trunc}{\Pi_{\mathsf{Trunc}}}
\newcommand{\MulTr}{\Pi_{\mathsf{MulTrunc}}}
\DeclareMathOperator{\negl}{negl}

\begin{document}

\title{Segmented Dealerless Function Secret Sharing\\ for Mixed-Width Private Inference}

\author{Anonymous Submission% <-- replace with author block for camera-ready
\thanks{Manuscript prepared from the architecture specification
\texttt{FSS\_Inference\_Architecture.md}. This draft targets a top-tier
security venue (NDSS/CCS/USENIX~Security/IEEE~S\&P) and is typeset in the
IEEE journal style for reference convenience.}}

\markboth{Preprint: Dealerless FSS Inference}%
{Author \MakeLowercase{\textit{et al.}}: Segmented Dealerless FSS for Mixed-Width Private Inference}

\maketitle

% =====================================================================
\begin{abstract}
Large neural networks commonly rely on wide fixed-point arithmetic to preserve
accuracy in private inference. Mixed precision reduces tensor and activation
interface widths, while signed rescaling of wider intermediates can still incur
high-bit-width nonlinear computation. Function Secret Sharing (FSS) supports
local nonlinear evaluation on masked inputs, but widely used inference systems
delegate key generation to a trusted third party. Existing full-tree distributed
generators remove this assumption at an exponential cost in the function's
input width, favoring short-domain nonlinear computation in practice. This
leaves a gap between the wide nonlinear domains used by accuracy-preserving
inference and the cost of generating their keys without a dealer.

We study a parallel multi-segment decomposition for dealerless FSS. The parties
split a public masked input into $\nseg$ digits, generate shallow segment trees,
and combine comparison results with a logarithmic-depth secure scan. For a
uniform split, tree generation from shared threshold bits requires
$\bigO(\nseg\,2^{\ell/\nseg})$ hash calls instead of $\bigO(2^\ell)$;
arithmetic-to-Boolean conversion and masked-input handling are accounted for
separately. We describe how these primitives compose with fixed-point inference
and distinguish tensor storage, accumulation, and nonlinear input widths.
Two-party localhost measurements demonstrate 32-bit segmented comparison
primitives. Cleartext ImageNette studies show the accuracy benefit of wider
fixed-point formats, including a 24-bit backbone with 16-bit ReLU interfaces.
A mixed-width FastGRNN execution retains 23- and 31-bit comparison domains in
signed rescaling, showing that narrow interfaces do not eliminate wide-domain
nonlinear work. The available evidence also includes a 30-sample SCI Google-30
consistency replay; an end-to-end comparison of SCI with the segmented backend
remains to be measured.
\end{abstract}

\begin{IEEEkeywords}
Function Secret Sharing, Distributed Comparison Function, Secure Inference,
Dealerless Key Generation, Correlated Oblivious Transfer, Secure Two-Party
Computation.
\end{IEEEkeywords}

% =====================================================================
\section{Introduction}\label{sec:intro}
\IEEEPARstart{S}{ecure} two-party computation (2PC) allows a model owner
and a data owner to evaluate a neural network while protecting model parameters
and inputs~\cite{secureml,cryptflow2}. Function Secret Sharing
(FSS)~\cite{bgi15,bgi16} is a useful building block for this task: a function
is split into two keys, and each party locally evaluates an output share on a
public masked input. FSS-based systems such as LLAMA~\cite{llama},
AriaNN~\cite{ariann}, Orca~\cite{orca}, and Sigma~\cite{sigma} exploit this
property to reduce online interaction for nonlinear operations.

Preserving accuracy in large neural networks commonly calls for wide fixed-point
arithmetic to accommodate both the range and fractional resolution of
intermediate values~\cite{cryptflow2,llama}. Mixed precision reduces widths where
the numerical program permits. SiRNN~\cite{sirnn}, for example, uses 8/16-bit
tensors and function interfaces alongside wider products and accumulations.
These formats require secure extension and rescaling between widths. Signed
rescaling can itself invoke high-bit-width comparisons, so narrowing tensor
storage or an activation interface does not remove all wide-domain nonlinear
computation.

Many FSS-based systems use a trusted dealer to generate keys and correlated masks.
This assumption matters when the two computing parties cannot delegate
preprocessing to a non-colluding third party. Distributed FSS instead generates
the keys jointly. Floram~\cite{floram} supports dealerless DPF generation with
Boolean output; Half-Tree~\cite{halftree} uses correlated GGM trees and parallel
COT; Distributed FSS~\cite{distfss} supports arithmetic-shared input and output.
These constructions provide the ingredients for a two-party inference path
without a trusted dealer, but their cost depends strongly on the domains of
the nonlinear operations.

For the full-tree distributed generators considered here, an $\ell$-bit domain
requires $\Theta(2^\ell)$ tree expansion work and sequential correction-word
exchanges across tree levels. At $\ell=32$, generating even the underlying
Half-Tree DPF requires approximately $1.5\cdot2^{32}$ hash calls. The
exponential generation cost makes these constructions more
practical for short-domain nonlinear operations and becomes a bottleneck when
inference repeatedly requires wide-domain comparisons, including those inside
mixed-precision conversions.

Our numerical and execution studies illustrate both sources of this demand.
On ImageNette, uniform 16-bit storage causes substantial accuracy losses in the
evaluated ResNets, while wider formats recover accuracy. Retaining a Q$(24,15)$
backbone also allows 16-bit ReLU interfaces with near-FP32 Top-1 accuracy.
Such interfaces still require conversion between numerical formats. In the
evaluated mixed-width FastGRNN execution, signed rescaling invokes 23- and
31-bit comparison domains despite the 8/16-bit tensor and activation interfaces
(Section~\ref{sec:eval}). We therefore target the high-bit-width nonlinear
operations that remain in accuracy-preserving fixed-point programs, including
mixed-precision programs. Selecting tensor widths is a separate numerical
optimization.

We address the tree-generation cost with a \emph{parallel multi-segment
decomposition}. An $\ell$-bit masked input is split into $\nseg$ public digits.
Each digit is evaluated against a shared threshold using a shallow tree, and a
balanced secure scan combines the results. For uniform segments, the exponential
term becomes $\nseg\,2^{\ell/\nseg}$. The same radix structure factors point
functions and comparisons. This mechanism applies at different function input
widths; selecting the numerical width of each tensor is a separate optimization.
Table~\ref{tab:related} places the construction among representative FSS systems.

Our contributions are as follows.
\begin{itemize}
  \item \textbf{A segmented dealerless comparison construction.}
  The radix identity in Lemma~\ref{lem:decomp} decomposes an $\ell$-bit
  comparison into reduced-domain comparisons and an equality scan.
  Theorem~\ref{thm:cost} characterizes tree generation from shared threshold
  bits and the subsequent scan. Arithmetic threshold decomposition and private
  query handling remain explicit parts of a complete comparison.

  \item \textbf{Composition with fixed-point nonlinear gates.}
  We combine Half-Tree correlated-GGM generation, distributed arithmetic
  interfaces, and Grotto-style prefix-parity evaluation. The resulting
  primitives support comparison, ReLU, reductions, and table-based activation
  constructions. We specify their arithmetic and communication interfaces
  separately from the numerical choices of a CNN or recurrent model.

  \item \textbf{Evidence for wide-domain nonlinear computation.}
  We report two-party localhost comparison measurements and cleartext
  fixed-point accuracy studies, together with an execution trace identifying
  wide comparison domains in mixed-precision rescaling. A historical Google-30
  replay checks SCI against SeeDot on 30 samples. These establish the starting
  point for a same-program comparison: the planned \emph{SCI + ours} backend retains
  SCI's linear layers, conversions, truncation, and activation approximations
  while replacing matching comparison/MSB primitives with segmented FSS.
\end{itemize}

\noindent\textbf{Organization.}
Section~\ref{sec:related} reviews related work. Section~\ref{sec:prelim} fixes
preliminaries. Section~\ref{sec:overview} gives the system overview and threat
model. Section~\ref{sec:core}, the technical core, develops the efficient
segmented dealerless FSS in two parts: the input decomposition
(\S\ref{sec:split}) and the Grotto-on-HalfTree segment generation
(\S\ref{sec:spine}). Section~\ref{sec:infer} composes these
primitives into inference gate constructions. Section~\ref{sec:eval} describes the
evaluation. Section~\ref{sec:conc} concludes.

% =====================================================================
\section{Related Work}\label{sec:related}

\subsection{FSS-based Secure Inference}
FSS reduces the online cost of comparison-type gates to a single round, which has
made it the backbone of several secure-inference systems. LLAMA~\cite{llama}
builds low-latency math (sigmoid, reciprocal) on DCF. AriaNN~\cite{ariann}
provides an FSS toolkit for private deep learning. Orca~\cite{orca} accelerates
FSS evaluation on GPUs. Pika~\cite{pika} evaluates arbitrary functions by
full-domain DPF lookups, and Grotto~\cite{grotto} introduces a parity-segment
tree that evaluates a DPF over an interval in $\bigO(\log N)$ rather than
$\bigO(N)$, scaling spline evaluation to $64$-bit words. All of these assume a
dealer (Grotto uses a non-colluding ``$+1$'' helper that supplies correlated
randomness offline), and Pika-style full-domain lookups inherit the
$\bigO(2^{n})$ cost that limits the practical domain. We borrow Grotto's
evaluation as one of our two evaluation paths, but run it on dealerlessly
generated keys over the shallow trees of our decomposition.

\subsection{Dealerless and Distributed FSS Key Generation}
Floram~\cite{floram} gives the first two-party DPF key generation but only for
boolean output. Half-Tree~\cite{halftree} introduces the correlated-GGM tree,
halving the cost of tree expansion in COT and DPF, and gives a distributed
DPF/DCF generation protocol from COT together with a parallel-COT preprocessing
subroutine, on both of which we build directly. Distributed FSS~\cite{distfss}
provides two-party DPF and DCF generation with both arithmetic-shared input and
output, resolves the arithmetic-comparison overflow problem, and characterizes the
$\bigO(\ell)$-round generation cost we attack, while explicitly naming
sublinear-round dealerless generation as open. Our decomposition is orthogonal to
these generators: it reduces the depth of the trees they expand. Their full-tree
expansion has exponential cost in the input width; Distributed FSS reports
arithmetic-sharing evaluations at 8, 16, and 18 bits. This cost limits the
practical use of repeated wide-domain nonlinear operations, including
conversions within mixed-precision inference. The 32-bit
primitive measurements in Section~\ref{sec:eval-split} evaluate this reduction.

\subsection{Secure Comparison Protocols}
Outside FSS, garbled circuits and the millionaires'-protocol line
(CrypTFlow2~\cite{cryptflow2}, SIRNN~\cite{sirnn}) realize comparison through
bit-decomposition and recursive digit comparison. Our radix decomposition
(Lemma~\ref{lem:decomp}) is the FSS analogue of CrypTFlow2's recursive
millionaire: it applies the same lexicographic identity, but to the
\emph{public} blinded value, so each digit comparison becomes an independent
reduced-domain FSS instance rather than a sequence of secure comparisons.

% --- prior-work matrix ---
\begin{table}[t]
\centering
\caption{Capabilities of representative FSS constructions.
\checkmark: supported, \ding{55}: not supported. The 32-bit column denotes
domain support, independently of key-generation cost or network precision.
ML target denotes intended application rather than completed end-to-end
evaluation. Our offline-OT property
applies to the preprocessed comparison core; the SIGMA activation path also
performs online preprocessing.}
\label{tab:related}
\renewcommand{\arraystretch}{1.2}
\setlength{\tabcolsep}{3pt}
\begin{tabular}{lccccc}
\toprule
System & Dealerless & Arith.\ I/O & $32$-bit & All-OT-offline & ML target\\
\midrule
LLAMA~\cite{llama}      & \ding{55} & \checkmark & \checkmark & \ding{55} & \checkmark\\
AriaNN~\cite{ariann}    & \ding{55} & \checkmark & \checkmark & \ding{55} & \checkmark\\
Grotto~\cite{grotto}    & $+1$ helper & \checkmark & \checkmark & partial & \checkmark\\
Pika~\cite{pika}        & \ding{55} & \checkmark & \ding{55} & \ding{55} & \checkmark\\
Floram~\cite{floram}    & \checkmark & \ding{55} & \checkmark & \ding{55} & \ding{55}\\
Half-Tree~\cite{halftree} & \checkmark & partial & \checkmark & \checkmark & \ding{55}\\
Dist.\ FSS~\cite{distfss} & \checkmark & \checkmark & \checkmark & \ding{55} & \ding{55}\\
\textbf{This work}      & \checkmark & \checkmark & \checkmark & partial & \checkmark\\
\bottomrule
\end{tabular}
\end{table}

% =====================================================================
\section{Preliminaries}\label{sec:prelim}

\subsection{Notation}
We describe arithmetic over $\Z_{2^{\ell}}$, embedding fixed-point reals
as $\lfloor 2^{s}\cdot z\rfloor \bmod 2^{\ell}$ for fractional scale $s$, and use
$\lambda=128$ as the computational security parameter. Table~\ref{tab:notation}
summarizes the notation used throughout the paper. A uniform split sets the segment
widths to $\wseg_j=\ell/\nseg$. Here $\ell$ denotes the domain of the
particular protocol instance; 32 bits is one measured configuration. A
mixed-width program separately specifies tensor storage and function I/O widths,
fractional scales, product and accumulator widths, and actual comparison or LUT
input widths. The implemented nonlinear operation determines the FSS domain.
In particular, signed rescaling can require a wide comparison even when tensor
storage and activation interfaces use fewer bits.

Following the cost convention of Distributed FSS~\cite{distfss}, communication
in analytical formulas is reported as the number of bits sent by one party;
for symmetric traffic the bidirectional total is twice that value. Empirical
measurements state their own counter boundary in Section~\ref{sec:eval}.
For a protocol $\Pi$ we write
$C^{\mathrm{off}}_{\Pi}$ and $R^{\mathrm{off}}_{\Pi}$ for the bits per party and
the round depth it spends before the input is known, and
$C^{\mathrm{on}}_{\Pi}$ and $R^{\mathrm{on}}_{\Pi}$ for the same two quantities
once evaluation begins. Parallel calls contribute the sum of their
communication but only the maximum of their round depths. Opening a $q$-bit
sharing therefore costs $q$ bits in one round. With OT-extension-based
triple generation~\cite{cryptflow2}, a Boolean \textsc{and} consumes one
preprocessed bit triple and costs $2$ online bits in one round; generating
its triple costs $\lambda+16$ offline bits. A conversion
$\mathsf{B2A}:\F_2\rightarrow\Z_{2^q}$ costs $\lambda+q$ bits in one round.
A $q$-bit Beaver triple costs $q(\lambda+16)$ offline bits to generate
in one round; consuming it opens two $q$-bit
differences, costing $2q$ bits in one round. Our $\BitBA$ preprocessing uses two such
multiplications to correlate its Boolean and arithmetic masks. Their triples
can be generated together, after which the two multiplications are dependent;
the resulting preprocessing costs $2q(\lambda+18)$ bits in three offline
rounds. Online, $\BitBA$ opens one masked ring element and one masked bit,
costing $q+1$ bits in one round.

\begin{table}[t]
\centering
\caption{Summary of notation.}
\label{tab:notation}
\renewcommand{\arraystretch}{1.2}
\begin{tabular}{ll}
\toprule
Notation & Description \\
\midrule
$P_0,P_1$ & the two parties \\
$b$ & party index, $b\in\{0,1\}$ \\
$\ell$ & bit-width of a protocol instance \\
$\lambda$ & computational security parameter, $\lambda=128$ \\
$\Z_{2^{\ell}}$ & arithmetic ring \\
$\F_2$ & binary field \\
$\R$ & ring-generic carrier, $\F_2$ or $\Z_{2^{\ell}}$ \\
$\ashare{x}$ & arithmetic share, $x=\ashare{x}_0+\ashare{x}_1$ \\
$\bshare{x}$ & boolean share, $x=\bshare{x}_0\oplus\bshare{x}_1$ \\
$\mathsf{B2A}$ & boolean-to-arithmetic conversion \\
$\BitBA$ & boolean-bit--arithmetic-value multiplication \\
$\nseg$ & segment count \\
$\wseg_j$ & width of segment $j$, $\sum_j\wseg_j=\ell$ \\
$C^{\mathrm{off}}_{\Pi},R^{\mathrm{off}}_{\Pi}$ & offline bits per party and rounds of $\Pi$ \\
$C^{\mathrm{on}}_{\Pi},R^{\mathrm{on}}_{\Pi}$ & online bits per party and rounds of $\Pi$ \\
$\xhat$ & masked input \\
$r_{\mathrm{in}},r_{\mathrm{out}}$ & input and output masks \\
$\Fcot$ & COT functionality \\
$\Delta$ & global offset \\
$K[r],M[r]$ & COT key and correlation, $M[r]=K[r]\oplus r\cdot\Delta$ \\
$s$ & GGM node label \\
$t$ & control  bit \\
$\mathit{CW}_i$ & public correction word at level $i$ \\
$\alpha$ & DPF special point \\
$\beta$ & DPF payload \\
$S$ & one-time hash key \\
$H_S(x)$ & keyed hash $H(S\oplus x)$, outputs $\lambda$ bits \\
$H'_S(x)$ & high $(\lambda{-}1)$-bit truncation of $H_S(x)$ \\
\bottomrule
\end{tabular}
\end{table}

\subsection{Cryptographic Primitives}
\textbf{Secret sharing}
Our work involves $2$-out-of-$2$ secret sharing over different rings. We describe the sharing semantics as follows:

\noindent\textbf{Arithmetic sharing $\ashare{x}$.} A secret value $x$ is additively shared between two parties such that $P_b$ holds $\ashare{x}_b$, satisfying $x=\ashare{x}_0+\ashare{x}_1 \bmod 2^{\ell}$.

\noindent\textbf{boolean sharing $\bshare{x}$.} A secret value $x$ is XOR-shared such that $P_b$ holds $\bshare{x}_b$, satisfying $x=\bshare{x}_0\oplus\bshare{x}_1$.

We write
$\mathsf{B2A}$ for boolean-to-arithmetic conversion. We use $\BitBA$ for
mixed-share multiplication: given a boolean sharing $\bshare{y}$ of a bit
$y\in\{0,1\}$ and an arithmetic sharing $\ashare{x}$ over $\Z_{2^\ell}$, it
returns an arithmetic sharing $\ashare{z}$ such that $z=x\cdot y$. Its
preprocessing prepares correlated masks for both sharing domains; online, the
parties open the masked $x$ and $y$ together in one round and derive
$\ashare{z}$ locally. This functionality is called $\mathsf{BitXA}$ by
Yang et al.~\cite{fssnn}; we write $\BitBA$ to make its boolean-bit and
arithmetic-value interface explicit.

\paragraph*{Ideal functionalities}
The protocols in this paper are analyzed in the $(\Fcot,\Frand)$-hybrid
model. $\Fcot$ is defined below. $\Frand$ is the standard coin-tossing
functionality: on invocation it outputs a uniform string $W\in\{0,1\}^{\lambda}$ to
both parties simultaneously, providing shared public randomness without either party
controlling the outcome.

\paragraph*{Arithmetic multiplication and truncation}
Beyond the mixed-share product $\BitBA$, recurrent and normalization layers
multiply two data-dependent arithmetic shares. We write $\Fmul$ for the
standard multiplication functionality over $\Z_{2^{\ell}}$: on
$\ashare{x},\ashare{y}$ it returns $\ashare{z}$ with $z=xy\bmod 2^{\ell}$,
extended componentwise to vectors, and realized from a multiplication triple
$(\ashare{a},\ashare{b},\ashare{ab})$ by the Beaver reduction: the parties
open $x-a$ and $y-b$ in one round and combine locally. Triples are the only
preprocessing $\Fmul$ needs, and they are generated dealerlessly from $\Fcot$
by the standard COT-to-triple construction: each of the $\ell$ bit positions of
each operand contributes one cross term, so one $\ell$-bit triple costs
$2\ell$ COTs and no trusted party. This is the same $\Fcot$ that $\PREP$
already batches, so $\Fmul$ adds no new setup assumption.

A fixed-point product carries $2f$ fractional bits and must be rescaled by $f$
before it is consumed downstream. We write $\Ftr$ for exact truncate-reduce:
on $\ashare{z}$ it returns an $\ashare{z'}$ in
$\mathbb Z_{2^{\ell-f}}$ such that
$z'=\lfloor z/2^f\rfloor\bmod 2^{\ell-f}$. Interpreting the result as a signed
two's-complement value gives the exact floor of the input's signed ring
representative divided by $2^f$. Section~\ref{sec:trunc} gives the dealerless
realization of SIGMA's low-bit DPF comparison for this functionality.

\paragraph*{The offset gate and masked evaluation}
\label{sec:offset}
For a gate $g$ we use the offset (masked) gate family
\begin{equation}
\hat g_{[r_{\mathrm{in}},r_{\mathrm{out}}]}(x) \;=\; g\!\left(x - r_{\mathrm{in}}\right) + r_{\mathrm{out}},
\label{eq:offset}
\end{equation}
where $r_{\mathrm{in}},r_{\mathrm{out}}$ are fresh random masks chosen at
generation time. The parties hold $\ashare{r_{\mathrm{in}}}$, locally form
$\ashare{\xhat}=\ashare{x}+\ashare{r_{\mathrm{in}}}$, and \emph{open} the masked
value $\xhat=x+r_{\mathrm{in}}\bmod 2^{\ell}$. Since $r_{\mathrm{in}}$ is uniform
and independent of $x$ and is never revealed, $\xhat$ is a one-time pad of $x$
and leaks nothing. Evaluating the FSS key (whose secret point is shifted by
$r_{\mathrm{in}}$ and whose output is shifted by $r_{\mathrm{out}}$) at the public
$\xhat$ yields additive shares of $g(x)$.


\textbf{Correlated OT and the Half-Tree}
\label{sec:cot}
A Correlated OT (COT) over $\F_{2^{\lambda}}$, modeled by $\Fcot$, gives a sender
$(\Delta, K[r])$ and a receiver $(r, M[r])$ with the COT correlation
\begin{equation}
M[r] \;=\; K[r]\oplus r\cdot\Delta, \qquad r\in\F_2,\ \Delta\in\F_{2^{\lambda}}.
\label{eq:cot}
\end{equation}
Following the COT-based multiplexer instantiation of~\cite{distfss}, a COT on an
$\ell$-bit value costs $\lambda{+}\ell$ bits of communication (an $\ell$-bit MUX
uses two such COTs at $2(\lambda{+}\ell)$ bits).
The half-tree (correlated-GGM, cGGM)~\cite{halftree} introduces a global offset
$\Delta$ and a single hash $H$, deriving the two children of a node $x$ as
$x_{\mathrm{left}}:=H(x)$ and $x_{\mathrm{right}}:=x-H(x)$, so that one child is
free and each level of the tree sums to $\Delta$, with
$H(x)=\pi(\sigma(x))\oplus\sigma(x)$ for a fixed-key permutation $\pi$ and a
linear orthomorphism $\sigma$. Because one child is free, a bare cGGM tree expands
in $\approx N$ random permutation calls, half the $\approx 2N$ a standard GGM tree
needs. The DPF tree we build keeps this one-call-per-node expansion on its first
$n-1$ levels but doubles the last level, where two $H$ calls per parent break the
leaf correlation so the leaves are pseudorandom. That last level adds $\approx N/2$
calls over the bare count, for a total of $\approx 1.5N$. The control bit at each
node coincides with the parity (control) bit that Grotto's evaluation reads, a
structural alignment we exploit in \S\ref{sec:spine}.

\subsection{Function Secret Sharing}\label{sec:fss}
For a function family $\mathcal{F}$, a function $f \in \mathcal{F}$ is split into two additive shares $f_0$ and $f_1$, such that for any $x$ in the domain of $f$, $f_0(x) + f_1(x) = f(x)$. While for $b \in \{0,1\}$, any single share $f_b$ does not reveal any information about $f$. The formal definition of FSS is as follows.

\paragraph{Definition 1} (Function Secret Sharing \cite{bgi15, bgi16}): Given the domain $\mathbb{G}_{in}$ and range $\mathbb{G}_{out}$, a (two-party) function secret sharing (FSS) scheme is a pair of probabilistic polynomial-time (PPT) algorithms (Gen, Eval) that satisfy:
\begin{itemize}
    \item $\text{Gen}(1^{\lambda},f)$: Given a security parameter $\lambda$ and a function $f \in \mathcal{F}$, the key generation algorithm generates two short keys $k_0, k_1$ that represent $f_0, f_1$, respectively, without revealing $f$.
    \item $\text{Eval}(b, k_b, x)$: Given the party index $b \in \{0, 1\}$, key $k_b$, and input $x \in \mathbb{G}_{in}$, the evaluation algorithm returns the additive share $f_b(x)$ such that $f_0(x) + f_1(x) = f(x)$.
\end{itemize}

\paragraph{Distributed Point Function (DPF) \cite{bgi16}} For $\alpha \in \mathbb{G}_{in}$ and $\beta \in \mathbb{G}_{out}$, the point function $f_{\alpha,\beta}$ is defined as $f_{\alpha,\beta}(\alpha) = \beta$, and $f_{\alpha,\beta}(x) = 0$ when $x \neq \alpha$. The two-party FSS construction of DPF has the following properties.
% detailed in Fig \ref{DPF1}

\textbf{DPF invariance.} Generation and evaluation of the FSS key are based on the GGM tree, where the nodes of the tree consist of a random label $s$ and a control bit $t$. The invariant property of DPF is that only the nodes along the $\alpha$ path have different labels with a control bit of 1, while the nodes on all other paths have identical labels with a control bit of 0.

\textbf{Correction words.} In the GGM tree, $s$ and $t$ are generated pseudo-randomly. To adhere to the DPF invariance, the dealer generates correction words for each layer along the $\alpha$ path to transform the random GGM-style tree. 

\textbf{Supporting arithmetic output.} The output of the basic DPF scheme only supports boolean sharing. To support the arithmetic output of an element $\beta \in \mathbb{G}_{out}$, Boyle et al. propose a masked transformation as Equation \ref{eq:DPF-last}. When $P_b$ holds control bit $t_b$ and leaf value $r_b$ on the last layer $\alpha$ path, the dealer computes the final correction word as $(-1)^{t_1} \cdot (\beta - r_0 - r_1)$, using $r_b$ to conceal $\beta$. $P_b$ locally computes the output by applying $t_b$ to the correction word. When $t_0 \neq t_1$, the output is $\beta$, otherwise, the output is 0.
\begin{equation}
  \label{eq:DPF-last}
  \begin{aligned} 
    (t_0 - t_1) \cdot (-1)^{t_1} \cdot (\beta - r_0 + r_1) + r_0 -r_1
  \end{aligned}
\end{equation}

\begin{figure}[h]
    \centering
   \includegraphics[width=\linewidth]{2PC.pdf} 
  \caption{\textbf{The traversal from the root to leaf in GGM tree with DPF invariant property on target $\alpha = 5$.}}
  \label{DPF1}
\end{figure}


% \subsection{Grotto's Parity-Segment Evaluation}
% \label{sec:grotto}
% Grotto~\cite{grotto} observes that a BGI shareholder can answer arbitrary segment
% (interval) queries on a DPF by running a \emph{prefix-parity} walk directly on its
% key share, at the cost of one half-PRG evaluation per traversed edge. Over a
% domain of size $N$ this evaluates a comparison or interval-containment test in
% $\bigO(\log N)$ rather than the $\bigO(N)$ of full-domain expansion, and the
% vector of segment parities is exactly the one-hot selection vector picking which
% piece of a piecewise function the secret input falls into. An inner product with
% deferred sign correction then lifts the result to the ring. Because this online
% phase is agnostic to how the key was produced, it composes with dealerless
% generation.

% =====================================================================
\section{System Overview and Threat Model}\label{sec:overview}

\subsection{Architecture}
The preprocessed comparison core runs in two phases (Fig.~\ref{fig:overview}). In the \emph{offline
phase}, $P_0$ and $P_1$ jointly run dealerless key generation: a single
parallel-COT preprocessing call ($\PREP$) produces a shared global offset
$\ashare{\Delta}$ and $\ell$ COT correlations on the secret threshold bits, from which the
parties generate, for every comparison gate, a set of $\nseg$ \emph{correlated
segment keys}, one shallow tree of depth $\wseg_j$ per segment, together with
the input mask $\ashare{r_{\mathrm{in}}}$. All OT communication occurs here, and
the transcript is independent of the actual inference inputs.

In the \emph{online phase}, for each comparison gate the parties (i) locally
blind and open $\xhat=x+r_{\mathrm{in}}$, (ii) locally split $\xhat$ into the
$\nseg$ public digits $\xhat\seg{1},\dots,\xhat\seg{\nseg}$, (iii) evaluate the
$\nseg$ segment keys \emph{in parallel}, and (iv) recombine the segment results
with a logarithmic-depth secure-\textsc{and} scan. The segment evaluation
uses boolean keys with a Grotto-style prefix-parity tree walk (\S\ref{sec:spine}).
For inference with private weights, linear layers also require secure
secret-by-secret multiplication and the prescribed rescaling. The planned
SCI + ours composition retains SCI's mixed-width linear and conversion paths.

\begin{figure*}[t]
\centering
\includegraphics[width=\textwidth]{figures/fig_overview.pdf}
\caption{Comparison architecture and intended gate composition, read left to
right. Two parties, $P_0$, the
\emph{user} holding private data $x$, and $P_1$, the \emph{model owner} holding a
commercial model $M$, compute with \emph{no trusted dealer}. In the
\textbf{offline} phase (all oblivious transfer, input-independent), a single
parallel-COT preprocessing call decomposes each depth-$\ell$ comparison tree into
$\nseg$ shallow segment trees, cutting dealerless generation from $\bigO(2^\ell)$
to $\bigO(\nseg\cdot 2^{\ell/\nseg})$ and yielding correlated segment keys
$\{k\seg{j}\}$. The keys hand off to the \textbf{online} phase (no OT): the parties
open the blinded value $\xhat=x+r_{\mathrm{in}}$, split it into $\nseg$ public
digits, evaluate the segment keys \emph{in parallel} via Grotto-style prefix-parity
tree walks, and stitch the results with a
logarithmic-depth secure-\textsc{and} scan to obtain the comparison share
$\ashare{\mathbf{1}\{x<r_{\mathrm{in}}\}}$. The right-hand panel denotes
inference applications of these gates; completed measurements are reported
separately in Section~\ref{sec:eval}.}
\label{fig:overview}
\end{figure*}

\subsection{Threat Model}
We adopt the standard simulation paradigm against a \emph{static, semi-honest}
PPT adversary that corrupts exactly one of $P_0,P_1$. The adversary follows the
protocol but tries to infer the honest party's input from its view. Security is
argued in the $(\Fcot,\Frand)$-hybrid model.

\begin{definition}[Security]
A protocol $\Pi$ securely realizes functionality $\mathcal G$ against a
semi-honest adversary corrupting $P_b$ if there is a PPT simulator
$\mathsf{Sim}_b$ such that
$\{\mathsf{Sim}_b(x_b,\mathcal G_b),\,\mathcal G\}\approx_c
\{\mathrm{view}^{\Pi}_b,\,\mathrm{out}^{\Pi}\}$,
where $\mathrm{view}^{\Pi}_b$ is $P_b$'s view and $\approx_c$ denotes
computational indistinguishability.
\end{definition}

Malicious security (input/key consistency, verifiable FSS) is out of scope and
deferred to future work, consistent with the dealerless-FSS line we build
on~\cite{distfss,halftree}.

% =====================================================================
\section{Segmented Dealerless FSS}\label{sec:core}
We first introduce input decomposition for a specified function domain
(\S\ref{sec:split}), then present dealerless segment key generation and
evaluation (\S\ref{sec:spine}).


\subsection{Parallel Decomposition}\label{sec:split}
Existing distributed FSS protocols~\cite{distfss,halftree} rely on a single 
depth-$\ell$ GGM tree, which incurs an infeasible $\Theta(2^{\ell})$ cost for 
large domains (e.g., $\ell=32$). To bypass this exponential bottleneck, 
we instead decompose the $\ell$-bit input into $\nseg$ aligned segments and 
evaluate each using a separate shallow tree.

\begin{figure*}[t]
  \centering
  \includegraphics[width=\textwidth]{figures/fig_split.pdf}
  \caption{\textbf{Decomposing the input into parallel segments.}
  (a) Overview of our $\nseg$-segment decomposition, which transforms an
  infeasible depth-$\ell$ tree into $\nseg$ parallel shallow trees.
  (b) A concrete instantiation of the composite comparison logic using
  $\xhat = 27|05|88|63$ and $r = 27|41|12|09$, demonstrating how the
  first diverging segment ($05 < 41$) strictly dictates the global order.
  Refer to this figure when reading Example~\ref{ex:split}, Definition~\ref{def:split},
  and Lemma~\ref{lem:decomp}.}
  \label{fig:split}
\end{figure*}

This decomposition is enabled by the offset gate from Eq.~\eqref{eq:offset}. 
Because the parties jointly open a blinded value $\xhat$ without revealing the 
underlying input $x$, we can directly partition this public $\xhat$ into $\nseg$ shorter segments.
Instead of traversing a single deep tree to compare $\xhat$ against 
a secret threshold $r$, we treat each segment as an independent 
FSS instance with its own GGM seed. The single deep tree therefore gives way 
to $\nseg$ shallow trees of depth $\wseg_j\approx\ell/\nseg$ (Fig.~\ref{fig:split}(a)).

Since these $\nseg$ trees are mutually independent, they can be generated and
evaluated in parallel. The computational cost is thus governed by the deepest sub-tree ($\max_j\wseg_j$) rather than the full depth ($\sum_j\wseg_j=\ell$). 
Every sub-tree then operates over a smaller domain of size $2^{\wseg_j}$, which is feasible for dealerless FSS generation and evaluation.
Building on this intuition, we formalize the segment-based encoding and its exact comparison logic below.

\textbf{The $\nseg$-Segment Encoding.}
Rather than treating $\xhat$ and $r$ as monolithic $\ell$-bit strings, we partition both into $\nseg$ aligned segments. A worked example shows why this suffices (Fig.~\ref{fig:split}b).\label{ex:split}
For $\ell=32$ bits split into $\nseg=4$ byte-wide segments, take
$\xhat = \mathtt{0x1B~05~58~3F}$ (decimal $27|5|88|63$) and
$r = \mathtt{0x1B~29~0C~09}$ (decimal $27|41|12|9$).
Scanning from the most significant byte, segment~1 matches ($27=27$), 
but segment~2 diverges ($5<41$), which already fixes $\xhat<r$. 
Segments~3 and~4 never need to be examined. 

We now state the segment encoding precisely.

\begin{definition}\label{def:split}
Let $\wseg_1,\dots,\wseg_{\nseg}$ be segment widths such that 
$\sum_{j=1}^{\nseg}\wseg_j=\ell$, indexed in decreasing order of significance. 
Let $W_j:=\sum_{i>j}\wseg_i$ denote the bit-offset of the $j$-th segment. 
For any $\xhat\in\Z_{2^{\ell}}$, its $j$-th segment $\xhat\seg{j}$ is defined as
\begin{equation}
\xhat\seg{j}:=\big\lfloor \xhat/2^{W_j}\big\rfloor \bmod 2^{\wseg_j},
\end{equation}
such that $\xhat=\sum_{j=1}^{\nseg}\xhat\seg{j}\cdot 2^{W_j}$.
\end{definition}

% The secret threshold $r \in \Z_{2^{\ell}}$ splits identically into segments $r\seg{j}$. Since $\xhat$ is public, the parties compute $\{\xhat\seg{j}\}_{j=1}^{\nseg}$ locally, but the secret segments $r\seg{j}$ stay embedded within the FSS keys, turning each segment into an independent FSS instance.
\begin{algorithm}[t]
\caption{Per-segment threshold decomposition ($\Pi_{\mathsf{Split}}$)}
\label{alg:split}
\textbf{Input:} $P_b$ holds the threshold share $\ashare{r}_b$ and the segment widths $(\wseg_1,\dots,\wseg_{\nseg})$. \\
\textbf{Output:} $P_b$ holds, for each segment $j$, the threshold-bit share $\bshare{r\seg{j}}_b$.

\vspace{0.3em}
\textbf{[Offline]} $P_0$ and $P_1$ jointly, with no dealer:
\begin{algorithmic}[1]
\State $\bshare{r}_b \gets \textsf{BitDec}(\ashare{r}_b)$ \Comment{the \textsf{BitDec} protocol of~\cite{distfss}}
\For{$j=1$ to $\nseg$}
  \State $\bshare{r\seg{j}}_b \gets \bshare{r}_b[W_j,\,W_j{+}\wseg_j)$ with bit-offset $W_j=\sum_{i>j}\wseg_i$
\EndFor
\State \textbf{return} $\big\{\bshare{r\seg{j}}_b\big\}_{j\in[\nseg]}$
\end{algorithmic}
\end{algorithm}

We realize this bit decomposition through the protocol $\Pi_{\mathsf{Split}}$
in Algorithm~\ref{alg:split}. The parties first run $\textsf{BitDec}$ once on the
$\ashare{r}_b$ to obtain boolean shares~\cite{distfss}, then slice out
the bits $\bshare{r\seg{j}}$ of each segment locally (Lines 2--3). Every segment
carries its own root seed into generation, so the $\nseg$
shallow GGM trees, each of depth $\wseg_j$, are built and evaluated concurrently,
which cuts the critical-path depth from $\ell$ down to $\max_j\wseg_j$. 

\textbf{Complexity.} The $\textsf{BitDec}$ protocol recovers $\bshare{r}$ from
$\ashare{r}$ via a two-party ripple-carry addition with carry recurrence
$c_{i+1}=g_i\oplus(p_i\wedge c_i)$, where $g_i=\ashare{r}_0[i]\wedge\ashare{r}_1[i]$
is the generate term and $p_i=\ashare{r}_0[i]\oplus\ashare{r}_1[i]$ is the propagate term.
The reference protocol of~\cite{distfss} spends three \textsc{and}s per bit
because it realizes $\mathcal{F}_{\mathsf{OR}}$ with a further \textsc{and}.
We observe that $g_i$ and $p_i\wedge c_i$ are mutually exclusive
($g_i=1$ forces $p_i=0$), so the disjunction reduces to the free \textsc{xor}
shown above, leaving $2\ell-3$ \textsc{and}s total: $\ell-1$ generate terms
plus $\ell-2$ carry gates.

Because the threshold $r$ is input-independent, the entire bit decomposition
executes offline. Segment slicing (Lines~2--3) is local and contributes no
communication. The $2\ell{-}3$ required Beaver triples cost
$(2\ell{-}3)(\lambda{+}16)$ bits to generate in one round. Consuming them
unfolds in $\ell{-}1$ rounds: the generate-layer \textsc{and}s ($\ell{-}1$ gates)
execute in parallel, followed by $\ell{-}2$ sequential carry \textsc{and}s, each
opening two masked bits for $4$ bits of communication. Offline,
$\Pi_{\mathsf{Split}}$ thus sends $(2\ell{-}3)(\lambda{+}20)$ bits per party across
$\ell$ rounds. Online, segment slicing is purely local, so $\Pi_{\mathsf{Split}}$
sends no bits and uses no rounds. The decomposition is omitted when the threshold
is sampled directly in Boolean shares, as in $\Pi_{\DReLU}$.

\textbf{Comparison decomposition.}
We formalize the above comparison logic over secret-shared segments as follows.

\begin{lemma}[Comparison decomposition]\label{lem:decomp}
For $\xhat,r\in\Z_{2^{\ell}}$ decomposed into segments as in Definition~\ref{def:split}, with
$L_j:=\ind{\xhat\seg{j}<r\seg{j}}$ and prefix-equality
$E_j:=\bigwedge_{i<j}\ind{\xhat\seg{i}=r\seg{i}}$ ($E_1=1$),
\begin{equation}
\ind{\xhat<r}\;=\;\bigoplus_{j=1}^{\nseg}\big(E_j\wedge L_j\big).
\label{eq:decomp}
\end{equation}
\end{lemma}
\begin{IEEEproof}
Let $j^\star$ be the index of the most-significant segment where $\xhat$ and $r$ differ. 
If $\xhat=r$, all $L_j=0$ and both sides evaluate to $0$. For $j<j^\star$, the segments are identical, 
yielding $L_j=0$. For $j>j^\star$, the prefix match is broken by segment $j^\star$, yielding $E_j=0$. At $j=j^\star$, 
the prefix matches ($E_{j^\star}=1$) and the term equals $L_{j^\star}=\ind{\xhat<r}$,
because the most-significant differing segment alone decides the order of the full values $\xhat$ and $r$.
Consequently, at most one summand is nonzero, allowing the $\bigoplus$ to coincide with logical disjunction.
\end{IEEEproof}

Since Lemma~\ref{lem:decomp} guarantees mutually exclusive non-zero summands,
the $\bigoplus$ reduces to a sum over the ring $\R$.
Denoting $e_i:=\ind{\xhat\seg{i}=r\seg{i}}$ as the per-segment equality bit,
we rewrite the comparison as an additive expansion gated by prefix products:
\begin{equation}\label{eq:decomp-ss}
\begin{aligned}
\rshare{\ind{\xhat<r}}
\;&=\;\sum_{j=1}^{\nseg}\Big(\rshare{L_j}\!\!\prod_{i=1}^{j-1}\rshare{e_i}\Big) \\
\;&=\;\rshare{L_1}\;+\;\rshare{e_1}\rshare{L_2}\;+\;\cdots \\
&\quad\;+\;\Big(\textstyle\prod_{i=1}^{\nseg-1}\rshare{e_i}\Big)\rshare{L_{\nseg}}.
\end{aligned}
\end{equation}

The same segment structure factorizes the point function as well, since equality on $\xhat$ holds exactly when it holds on every segment:
\begin{equation}
\rshare{\ind{\xhat=\alpha}}\;=\;\prod_{j=1}^{\nseg}\rshare{\ind{\xhat\seg{j}=\alpha\seg{j}}}
\;=\;\prod_{j=1}^{\nseg}\rshare{e_j}.
\label{eq:dpf-decomp}
\end{equation}

Finally, we combine the pairs $(L_j,e_j)$ with the associative operator
$(L,e)\otimes(L',e')=(L\oplus(e\wedge L'),e\wedge e')$. A balanced reduction
uses exactly $2(\nseg-1)$ Boolean \textsc{and} gates in
$\lceil\log_2\nseg\rceil$ dependent rounds. Under the convention above, its
online communication is $4(\nseg-1)$ bits per party.


\begin{theorem}[Generation and evaluation cost]\label{thm:cost}
Let $d=\max_j\wseg_j$. Dealerless generation of the $\nseg$ segment keys from
the Boolean shares of an $\ell$-bit threshold costs
\begin{equation}
\begin{aligned}
R_{\mathrm{key}}&=d+2,\\
C_{\mathrm{key}}(\ell,\boldsymbol{\wseg})
&=\lambda+2\ell+1+
\sum_{j=1}^{\nseg}\big((\wseg_j-1)(\lambda+1)+\lambda+2\big)
\end{aligned}
\label{eq:key-cost}
\end{equation}
offline rounds and bits per party, respectively, and performs
$H_{\mathrm{key}}=\sum_j(3\cdot2^{\wseg_j-1}+3)$ hash calls per party. The
$3\cdot2^{\wseg_j-1}$ term consists of one hash for every non-leaf node and
two hashes for every leaf parent; the four leaf-mask hashes yield the constant
$+3$ after accounting for the root offset. Preprocessing the
$2(\nseg-1)$ bit triples required by the cross-segment reduction adds
$4(\nseg-1)(\lambda+1)$ offline bits and can be batched with $\PREP$, so it
does not increase $R_{\mathrm{key}}$. Once the query is public, evaluation
costs $\lceil\log_2\nseg\rceil$ online rounds and $4(\nseg-1)$ bits per party;
opening an $\ell$-bit masked query adds one round and $\ell$ bits.
\end{theorem}
\begin{IEEEproof}
$\PREP$ takes two rounds and $\lambda+2\ell+1$ bits per party. For segment
$j$, its $\wseg_j-1$ non-leaf correction words cost
$(\wseg_j-1)(\lambda+1)$ bits and its leaf correction word costs
$\lambda+2$ bits. Equal-depth layers of all segments are sent together, hence
the $d+2$ round depth and the sum in~\eqref{eq:key-cost}. A depth-$\wseg_j$
tree has $2^{\wseg_j-1}-1$ non-leaf nodes; its leaf stage makes two hash calls
per leaf parent plus four mask hashes. The balanced operator above has
$\nseg-1$ internal nodes and performs two \textsc{and}s at each node, which
gives the stated triple count, communication, and round depth.
\end{IEEEproof}

For a uniform split $\wseg_j=\ell/\nseg$, Equation~\eqref{eq:key-cost}
simplifies to
$C_{\mathrm{key}}=(\ell+1)\lambda+3\ell+\nseg+1$ bits per party,
$R_{\mathrm{key}}=\ell/\nseg+2$, and
$H_{\mathrm{key}}=\nseg(3\cdot2^{\ell/\nseg-1}+3)$. Thus the computation remains
$\bigO(\nseg2^{\ell/\nseg})$, while the complete masked comparison takes
$\ell/\nseg+3+\lceil\log_2\nseg\rceil$ offline-plus-online rounds when the two
phases are counted consecutively.
This bound starts from shared threshold bits. Converting an arithmetic-shared
threshold with Algorithm~\ref{alg:split} adds its own preprocessing cost.

\textbf{Efficiency Gains and Parameter Selection.} The exponential reduction of 
Theorem~\ref{thm:cost} has diminishing returns: the term $C(\nseg)=\nseg\,2^{\ell/\nseg}$ decays 
rapidly in $\nseg$, while the linear overhead of maintaining $\nseg$ separate trees grows 
monotonically. The optimum therefore settles at a small $\nseg$, trading round complexity 
against bandwidth in a way that depends on the network. We treat $\nseg$ as a free parameter 
and determine its optimal value experimentally in Section~\ref{sec:eval}.


\subsection{Segment Key Generation and Evaluation}\label{sec:spine}
GGM tree expansion constitutes the core mechanism for instantiating FSS primitives \cite{bgi15,bgi16}.
As established in Section~\ref{sec:fss}, generating a DPF divides into two parts: the non-leaf
layers ($1$ to $d{-}1$) expand the intermediate nodes and compute one correction word per layer
to preserve the DPF invariant, while the leaf layer converts the final nodes into secret shares
over the output range. We generate both parts for each shallow segment tree without a dealer.

We follow the core insight of Grotto~\cite{grotto} that
a single DPF suffices for comparison. The prefix parity at the
evaluated index already yields the comparison bit $\ind{\xhat\seg{j}<r\seg{j}}$.
Therefore, generating one point function per segment covers
the nonlinear computation of PPML and avoids maintaining a
separate comparison key alongside each equality key.
To generate these DPFs without a dealer, we adapt existing
primitives into a composite construction. We employ the correlated GGM (cGGM)
expansion of Half-tree~\cite{halftree} to grow the shallow segment
trees from a single parallelized offline COT invocation, which halves
the per-node hashing cost and the communication overhead. The construction proceeds in
three stages: (1) parallel-COT preprocessing ($\PREP$, Algorithm~\ref{alg:prep}),
which resolves every dependency on the secret threshold; (2) cGGM expansion of layers
$1$ to $d{-}1$; and (3) boolean leaf recovery at layer $d$. Stages (2) and (3) form the
single generation procedure $\Pi_{\mathsf{Seg}}$ (Algorithm~\ref{alg:gen}). All stages are
identical across the $\nseg$ segment trees.

% \begin{figure*}[t]
%   \centering
%   \includegraphics[width=\textwidth]{figures/fig_compose.pdf}
%   \caption{\textbf{Segment key generation architecture.}
%   Dealerless offline preprocessing: parallel-COT ($\PREP$) followed by
%   $\Pi_{\mathsf{Seg}}$ (Algorithm~\ref{alg:gen}), which expands layers $1 \dots d{-}1$
%   via cGGM and recovers the boolean leaf at layer $d$ in 1 round. Online,
%   Grotto-style prefix-parity tree walk yields the boolean comparison bit, which
%   is retained for the multiplication-scan of Section~\ref{sec:split} over $\F_2$
%   or combined with an arithmetic payload through $\BitBA$.}
%   \label{fig:compose}
% \end{figure*}


\begin{algorithm}[t]
\caption{Parallel-COT preprocessing ($\PREP$)}
\label{alg:prep}
\textbf{Input:} $P_b$ (for $b \in \{0,1\}$) holds the secret shares $\bshare{\alpha_i}_b$ for $i\in[1,\ell]$ and a global key $\Delta'_b\in\{0,1\}^{\lambda}$. \\
\textbf{Output:} $P_b$ obtains a share $\bshare{\Delta}_b$ with
$\lsb(\bshare{\Delta}_0\oplus\bshare{\Delta}_1)=1$ and COT correlations
$\{(K[\bshare{\alpha_i}_{1-b}],\,M[\bshare{\alpha_i}_b])\}_{i\in[\ell]}$ such that 
$M[\bshare{\alpha_i}_b]=K[\bshare{\alpha_i}_b]\oplus\bshare{\alpha_i}_b\cdot\bshare{\Delta}_{1-b}$.

\vspace{0.3em}
\textbf{[Offline]} $P_0$ and $P_1$ jointly, with no dealer:
\begin{algorithmic}[1]
\State The parties exchange $\lsb(\Delta'_b)$ and $P_b$ locally sets
       $\bshare{\Delta}_b:=\Delta'_b\oplus(0^{\lambda-1}\,\|\,(\lsb(\Delta'_{1-b})\oplus b))$ 
       so the two local keys yield one shared $\Delta$ with $\lsb(\Delta)=1$. 
\State The parties invoke the $\Fcot$ functionality to generate $\ell$ correlations 
in parallel. As a receiver, $P_b$ obtains choice bits $r^i_b$ and values $M'[r^i_b]$ 
paired with keys $K'[r^i_b]$ held by the sender $P_{1-b}$, so that 
$M'[r^i_b]=K'[r^i_b]\oplus r^i_b\cdot\Delta'_{1-b}$ for $i\in[1,\ell]$. 
\State $P_b$ computes $g^i_b := \bshare{\alpha_i}_b \oplus r^i_b$ for $i\in[1,\ell]$ and sends $(g^1_b,\dots,g^{\ell}_b)$ to $P_{1-b}$.
       \Comment{retarget the random COTs onto $\bshare{\alpha_i}$}
\For{$i=1$ to $\ell$} \Comment{$P_b$ updates correlation locally}
  \State $K[\bshare{\alpha_i}_{1-b}] := K'[r^i_{1-b}] \oplus (g^i_{1-b} \cdot \bshare{\Delta}_b)$.
  \State $M[\bshare{\alpha_i}_b] := M'[r^i_b] \oplus \big(r^i_b \cdot (0^{\lambda-1} \| (\lsb(\Delta'_{1-b}) \oplus b))\big)$.
\EndFor
\State \textbf{return} $\big(\bshare{\Delta}_b,\,\{(K[\bshare{\alpha_i}_{1-b}],\,M[\bshare{\alpha_i}_b])\}_{i\in[1,\ell]}\big)$
\end{algorithmic}
\end{algorithm}

Following Half-tree \cite{halftree}, our segment-key generation initializes with a parallel-COT preprocessing phase
(Algorithm~\ref{alg:prep}). The parties combine their local keys $\Delta'_b$
into a shared offset $\bshare{\Delta}$ with $\lsb(\Delta)=1$. They then draw $\ell$ random COT correlations via a
single parallel $\Fcot$ invocation and retarget them to $\alpha_i$ using one round of difference bits $g^i_b$ (Lines 2--4).
Consequently, all cryptographic dependencies on the secret $\alpha$ are resolved in exactly two rounds:
one parallel COT invocation and one retargeting round that exchanges both $\lsb(\Delta'_b)$ and the difference vector, irrespective of the number of instantiated trees.

Building on this preprocessing, $\Pi_{\mathsf{Seg}}$ (Algorithm~\ref{alg:gen})
generates each segment tree in full. The non-leaf stage (Lines 1--5) expands layers
$1$ to $d{-}1$: to guarantee mutually independent root shares across the $\nseg$ trees
while safely reusing the shared offset, the protocol re-randomizes the root of segment
$j$ by XORing $\bshare{\Delta}_b$ with a fresh public coin $W$ (Line 1). Each subsequent
layer reveals the public correction word $\mathit{CW}_i$ by combining the hashes $H_S$
and the $\PREP$ correlations (Line 3), securely blinding the $\alpha$ path. Because the
cGGM derives both children of a node from a single hash invocation (the left
child directly and the right as its complement), the per-level hash cost drops
to one per node instead of two (Line 4).

The leaf stage (Lines 6--9) recovers layer $d$. Lifting the leaves to
arithmetic shares offline would increase dealerless preprocessing. We keep
the leaves boolean and defer the arithmetic
conversion to the online phase. $P_b$ folds the depth-$(d{-}1)$ nodes into the cross-hash
arrays $XH_b[\sigma],XL_b[\sigma]$ (Line 6), samples a mask $\mu_b$ and reshares the
correction-bit material (Line 7), then reconstructs the boolean leaf correction word
$\mathit{CW}_d$ in one round of $\lambda{+}2$ bits (Lines 8--9). The resulting key drives
a purely local tree walk online, whose control bit at the evaluated index is the boolean
comparison share $\bshare{u}_b$; one $\BitBA$ call directly multiplies this bit
by the arithmetic payload and returns the arithmetic share $\ashare{w}_b$.

% \begin{algorithm}[t]
% \caption{Per-segment key generation ($\Pi_{\mathsf{Seg}}$)}
% \label{alg:gen}
% \textbf{Input:} $P_b$ holds the segment-bit shares $\bshare{\alpha\seg{j}}_b$ (from $\Pi_{\mathsf{Split}}$), the offset share $\bshare{\Delta}_b$ and correlations $\{(K[\cdot],M[\cdot])\}$ (from $\PREP$), a one-time hash key $S$, and the signed payload $\ashare{\beta}_b$, with segment depth $d=\wseg_j$. \\
% \textbf{Output:} $P_b$ holds the boolean key $k_b$, and online the arithmetic share $\ashare{w}_b$.

% \vspace{0.3em}
% \textbf{[Offline]} $P_0$ and $P_1$ jointly, with no dealer:
% \begin{algorithmic}[1]
% \State Sample a public coin $W \leftarrow \Frand$ and set the root $(s^{0,0}_b\|t^{0,0}_b):=\bshare{\Delta}_b\oplus W$.
% \For{$i=1$ to $d-1$} \Comment{non-leaf layers}
% \State $\bshare{\mathit{CW}_i}_b := \big(\bigoplus_{j \in [0,2^{i-1})} H_S(s^{i-1,j}_b\|t^{i-1,j}_b)\big) \oplus \overline{\bshare{\alpha_i}_b} \!\cdot\!\bshare{\Delta}_b\oplus K[\bshare{\alpha_i}_{1-b}]\oplus M[\bshare{\alpha_i}_b]$
% and $P_b$ sends it to $P_{1-b}$ to reconstruct $\mathit{CW}_i$.
% \State $(s^{i,2j+\sigma}_b\|t^{i,2j+\sigma}_b) :=
% H_S(s^{i-1,j}_b\|t^{i-1,j}_b)\;\oplus\;\sigma\cdot(s^{i-1,j}_b\|t^{i-1,j}_b)\;\oplus\;t^{i-1,j}_b
% \cdot\mathit{CW}_i$ for $j\in[0,2^{i-1}),\sigma\in\{0,1\}$
% \EndFor
% \State $(XH_b[\sigma]\|XL_b[\sigma]):=\bigoplus_{p \in [0,2^{d-1})} H_S\big((s^{d-1,p}_b\|t^{d-1,p}_b)\oplus\sigma\big)$ for $\sigma\in\{0,1\}$ \Comment{$XH_b$: high $\lambda-1$ bits, $XL_b$: last bit}
% \State $P_b$ samples $\mu_b$, sets $d_b:=H'_S(\mu_b\oplus K[\bshare{\alpha_d}_{1-b}])\oplus H'_S(\mu_b\oplus K[\bshare{\alpha_d}_{1-b}]\oplus\bshare{\Delta}_b)\oplus XH_b[0]\oplus XH_b[1]$, and reshares $(\mu_b,d_b)$
% \State $P_b$ forms the share
% \Statex $\quad\begin{aligned}
% \bshare{HCW}_b &:= XH_b[\overline{\bshare{\alpha_d}_b}]\oplus H'_S(\mu_b\oplus K[\bshare{\alpha_d}_{1-b}]) \\
% &\quad \oplus\, H'_S(\mu_{1-b}\oplus M[\bshare{\alpha_d}_b])\oplus \bshare{\alpha_d}_b\cdot d_{1-b},
% \end{aligned}$
% \Statex $\quad \bshare{LCW[\sigma]}_b := XL_b[\sigma]\oplus\bshare{\alpha_d}_b\oplus b\{\sigma{=}0\}$ for $\sigma\in\{0,1\}$
% \State Reconstruct $\mathit{CW}_d:=(HCW\|LCW[0]\|LCW[1])$ and \textbf{return} the boolean key $k_b=\big((s^{0,0}_b\|t^{0,0}_b),\{\mathit{CW}_i\}_{i\le d}\big)$
% \end{algorithmic}

% \vspace{0.3em}
% \textbf{[Online]} On the opened $\hat{x}$, each $P_b$ recovers the arithmetic share:
% \begin{algorithmic}[1]
% \State Parse $k_b=\big((s^{0,0}_b\|t^{0,0}_b),\{\mathit{CW}_i\}_{i\le d}\big)$ and extract the segment digit $\xhat\seg{j}=(\xhat\seg{j}_1,\dots,\xhat\seg{j}_d)$.
% \State $\bshare{L_j}_b\gets 0$ \Comment{initialize the comparison accumulator}
% \For{$i=1$ to $d$} 
% \State $(s^L_b\|t^L_b)\|(s^R_b\|t^R_b) := H_S(s^{i-1}_b) \oplus t^{i-1}_b\cdot\mathit{CW}_i$ 
% \If{$\xhat\seg{j}_i = 0$}
% \State $\bshare{L_j}_b \gets \bshare{L_j}_b \oplus t^R_b$; $(s^i_b\|t^i_b)\gets(s^L_b\|t^L_b)$ 
% \Else
% \State $(s^i_b\|t^i_b)\gets(s^R_b\|t^R_b)$ 
% \EndIf
% \EndFor
% \State $\bshare{e_j}_b \gets t^d_b$ \Comment{leaf control bit is the equality share}
% \State \textbf{return} $\ashare{w}_b\gets\BitBA(\bshare{L_j},\ashare{\beta})_b$.
% \end{algorithmic}
% \end{algorithm}

\begin{algorithm}[t]
\caption{Per-segment key generation and evaluation ($\Pi_{\mathsf{Seg}}$)}
\label{alg:gen}
\textbf{Input:} $P_b$ holds the segment-bit shares $\bshare{\alpha\seg{j}}_b$ (from $\Pi_{\mathsf{Split}}$), the offset share $\bshare{\Delta}_b$ and correlations $\{(K[\cdot],M[\cdot])\}$ (from $\PREP$), a one-time hash key $S$, and the signed payload $\ashare{\beta}_b$, with segment depth $d=\wseg_j$. \\
\textbf{Output:} $P_b$ holds the boolean key $k_b$, and online the arithmetic share $\ashare{w}_b$.

\vspace{0.3em}
\textbf{[Offline]} $P_0$ and $P_1$ jointly, with no dealer:
\begin{algorithmic}[1]
\State Sample a public coin $W \leftarrow \Frand$ and set the root $(s^{0,0}_b\|t^{0,0}_b):=\bshare{\Delta}_b\oplus W$.
\For{$i=1$ to $d-1$} 
\State $\bshare{\mathit{CW}_i}_b := \big(\bigoplus_{j \in [0,2^{i-1})} H_S(s^{i-1,j}_b\|t^{i-1,j}_b)\big) \oplus \overline{\bshare{\alpha_i}_b} \!\cdot\!\bshare{\Delta}_b\oplus K[\bshare{\alpha_i}_{1-b}]\oplus M[\bshare{\alpha_i}_b]$
and $P_b$ sends it to $P_{1-b}$ to reconstruct $\mathit{CW}_i$.
\State $(s^{i,2j+\sigma}_b\|t^{i,2j+\sigma}_b) :=
H_S(s^{i-1,j}_b\|t^{i-1,j}_b)\;\oplus\;\sigma\cdot(s^{i-1,j}_b\|t^{i-1,j}_b)\;\oplus\;t^{i-1,j}_b
\cdot\mathit{CW}_i$ for $j\in[0,2^{i-1}),\sigma\in\{0,1\}$
\EndFor
\State $(XH_b[\sigma]\|XL_b[\sigma]):=\bigoplus_{p \in [0,2^{d-1})} H_S\big((s^{d-1,p}_b\|t^{d-1,p}_b)\oplus\sigma\big)$ for $\sigma\in\{0,1\}$ \Comment{$XH_b$: high $\lambda-1$ bits, $XL_b$: last bit}
\State $P_b$ samples $\mu_b$, sets $d_b:=H'_S(\mu_b\oplus K[\bshare{\alpha_d}_{1-b}])\oplus H'_S(\mu_b\oplus K[\bshare{\alpha_d}_{1-b}]\oplus\bshare{\Delta}_b)\oplus XH_b[0]\oplus XH_b[1]$, and reshares $(\mu_b,d_b)$
\State $P_b$ forms the share
\Statex $\quad\begin{aligned}
\bshare{HCW}_b &:= XH_b[\overline{\bshare{\alpha_d}_b}]\oplus H'_S(\mu_b\oplus K[\bshare{\alpha_d}_{1-b}]) \\
&\quad \oplus\, H'_S(\mu_{1-b}\oplus M[\bshare{\alpha_d}_b])\oplus \bshare{\alpha_d}_b\cdot d_{1-b},
\end{aligned}$
\Statex $\quad \bshare{LCW[\sigma]}_b := XL_b[\sigma]\oplus\bshare{\alpha_d}_b\oplus b\{\sigma{=}0\}$ for $\sigma\in\{0,1\}$
\State Reconstruct $\mathit{CW}_d:=(HCW\|LCW[0]\|LCW[1])$ and \textbf{return} the boolean key $k_b=\big((s^{0,0}_b\|t^{0,0}_b),\{\mathit{CW}_i\}_{i\le d}\big)$
\end{algorithmic}

\vspace{0.3em}
\textbf{[Online]} On the opened $\hat{x}$, each $P_b$ recovers the arithmetic share:
\begin{algorithmic}[1]
\State Parse $k_b=\big((s^{0,0}_b\|t^{0,0}_b),\{\mathit{CW}_i\}_{i\le d}\big)$ and extract the segment digit $\xhat\seg{j}=(\xhat\seg{j}_1,\dots,\xhat\seg{j}_d)$.
\State $\bshare{L_j}_b\gets 0$ \Comment{initialize the comparison accumulator}
\For{$i=1$ to $d$} 
\State $(s^L_b\|t^L_b)\|(s^R_b\|t^R_b) := H_S(s^{i-1}_b) \oplus t^{i-1}_b\cdot\mathit{CW}_i$ 
\If{$\xhat\seg{j}_i = 0$}
\State $\bshare{L_j}_b \gets \bshare{L_j}_b \oplus t^R_b$; $(s^i_b\|t^i_b)\gets(s^L_b\|t^L_b)$ 
\Else
\State $(s^i_b\|t^i_b)\gets(s^R_b\|t^R_b)$ 
\EndIf
\EndFor
\State $\bshare{e_j}_b \gets t^d_b$ \Comment{leaf control bit is the equality share}
\State \textbf{return} $\ashare{w}_b\gets\BitBA(\bshare{L_j},\ashare{\beta})_b$.
\end{algorithmic}
\end{algorithm}

The online phase (Lines 1--11 of the \textbf{[Online]} block) implements the
Grotto prefix-parity mechanism. Each party initializes a running accumulator
$\bshare{L_j}_b$ to zero (Line 2), then walks the tree from root to the leaf
indexed by $\xhat\seg{j}$. At each level $i$, the party expands both children
from the current node with the correction word applied (Line 4), descends along
the path specified by bit $\xhat\seg{j}_i$, and crucially, XORs the right child's
control bit $t^R_b$ into $\bshare{L_j}_b$ whenever $\xhat\seg{j}_i=0$ (Lines 5--6).
This accumulated XOR of off-path sibling tags yields the boolean share
$\bshare{L_j}_b$ of the comparison $\ind{\xhat\seg{j}<r\seg{j}}$ after $d$ layers,
while the leaf control bit $t^d_b$ separately captures the equality share
$\bshare{e_j}_b$ (Line 10). The walk touches only $d$ nodes and requires no
communication. A single $\BitBA$ call directly multiplies $\bshare{L_j}$ by
the arithmetic payload $\ashare{\beta}$ and returns $\ashare{w}_b$ (Line 11),
the only online round that consumes preprocessed correlations. The boolean
comparison and equality outputs feed the multiplication-scan step
(\S\ref{sec:split}) over $\F_2$. Appendix~\ref{app:full} states how the
per-segment comparison and equality bits combine across the $\nseg$ segments.

\textbf{Complexity.} For all $\nseg$ trees together, $\PREP$ costs $2$ offline
rounds and $\lambda+2\ell+1$ bits per party. The non-leaf and leaf correction
words add
$\sum_j((\wseg_j-1)(\lambda+1)+\lambda+2)$ bits and
$\max_j\wseg_j$ rounds because equal-depth words are batched across segments.
Thus key generation has exactly the cost $R_{\mathrm{key}}$ and
$C_{\mathrm{key}}$ in~\eqref{eq:key-cost}, together with
$\sum_j(3\cdot2^{\wseg_j-1}+3)$ hash calls. Evaluation at one public digit per
segment is local: it uses $\sum_j(\wseg_j+1)=\ell+\nseg$ on-path hashes and no
communication.
If an arithmetic payload of $p$ bits is requested for every segment, the
$\nseg$ parallel $\BitBA$ calls add $4\nseg p(\lambda+p+1)$ offline bits in
three preprocessing rounds and $\nseg(p+1)$ online bits in one round. The
comparison protocol retains the Boolean $(L_j,e_j)$ outputs instead and pays
only the cross-segment reduction stated in Theorem~\ref{thm:cost}.
Thus the arithmetic-payload variant has the exact totals
\begin{equation}
\begin{aligned}
C^{\mathrm{off}}_{\mathsf{Seg},p}
 &=C_{\mathrm{key}}(\ell,\boldsymbol{\wseg})
   +4\nseg p(\lambda+p+1),\\
R^{\mathrm{off}}_{\mathsf{Seg},p}&=\max\{d+2,3\},\\
C^{\mathrm{on}}_{\mathsf{Seg},p}&=\nseg(p+1),\qquad
R^{\mathrm{on}}_{\mathsf{Seg},p}=1.
\end{aligned}
\label{eq:segment-payload-cost}
\end{equation}

% =====================================================================


\section{Nonlinear Gates and Inference Composition}\label{sec:infer}
Built upon the dealerless FSS key generation of Section~\ref{sec:core}, this
section provides the secure protocols that a neural-network inference requires,
including DReLU, the lookup table, ReLU, MaxPool, and ArgMax. The gate
constructions and recurrent-cell cost formulas specify a composition; the
measurements in Section~\ref{sec:eval} distinguish implemented primitives from
complete secure network execution.
% We start with DReLU (\S\ref{sec:Drelu}), the sign-extraction gate that underlies every signed comparison, and the lookup table (\S\ref{sec:lut}), which evaluates nonlinear activations exactly rather than through spline approximation. From these two primitives we derive ReLU, MaxPool, and ArgMax through standard reductions (\S\ref{sec:derived}). We then assemble the components into the inference protocol $\Pi_{\mathsf{Infer}}$ (\S\ref{sec:infer-proto}) and analyze its cost and security. 
% Each protocol places all data-independent work in the offline phase and keeps a single opening as the only necessary online interaction per gate.

\subsection{DReLU}\label{sec:Drelu}
DReLU extracts the sign of a shared signed integer by computing its most significant
bit (MSB). For a signed integer $x \in \Z_{2^\ell}$ with arithmetic sharing
$x=\ashare{x}_0+\ashare{x}_1$, the protocol takes $\ashare{x}_b$ as input and
outputs the boolean share
$\bshare{\DReLU(x)}_b=\bshare{\ind{x<2^{\ell-1}}}_b
=\bshare{1-\mathrm{MSB}(x)}_b$.

Extracting the MSB from an arithmetic sum over $\Z_{2^\ell}$ requires a carry correction. 
For share values $x_0,x_1$ with $x=(x_0+x_1)\bmod 2^\ell$, 
standard approaches~\cite{cryptflow2,sigma,falcon} rely on the following identity:
\begin{equation}
  \label{eq:MSB}
  \mathrm{MSB}(x)=\mathrm{MSB}(x_0)\oplus\mathrm{MSB}(x_1)
  \oplus\ind{y_0+y_1\ge 2^{\ell-1}}.
\end{equation}
where $y_b:=x_b\bmod 2^{\ell-1}$.

To avoid expensive online bit-decomposition, we apply this 
carry logic to a masked input $\hat{x}=x+r\bmod 2^\ell$. 
Let $r_H:=r_{\ell-1}$, $r_L:=r\bmod 2^{\ell-1}$, and 
$\hat{y}:=\hat{x}\bmod 2^{\ell-1}$. Since $x=\hat{x}-r\bmod 2^\ell$, 
subtracting the lower $\ell-1$ bits produces a borrow exactly 
when $\hat{y}<r_L$, yielding:
\begin{equation}
\label{eq:MSB_last}
\mathrm{MSB}(x)=\mathrm{MSB}(\hat{x})\oplus r_H\oplus\ind{\hat{y}<r_L}.
\end{equation}
We substitute this result into the 
identity $\DReLU(x)=1-\mathrm{MSB}(x)$ to obtain our target evaluation function:
\begin{equation}
  \label{eq:drelu-fold}
  \DReLU(x)=\underbrace{(r_H\oplus1)}_{\text{offline}}
  \oplus \underbrace{\ind{\hat{y}<r_L}}_{\text{online: }\mathit{res}} 
  \oplus \underbrace{\mathrm{MSB}(\hat{x})}_{\text{online: public}}.
\end{equation}

We present the detailed protocol $\Pi_{\DReLU}$ in Algorithm~\ref{alg:drelu}.
During the offline phase, the parties prepare the mask $r$. Each party $P_b$ 
samples independent random bits $\rho_{i,b} \in \{0,1\}$ to form boolean 
shares $\bshare{r_i}_b:=\rho_{i,b}$, ensuring each 
bit $r_i=\rho_{i,0}\oplus\rho_{i,1}$ is uniformly random. 
The parties then perform $\ell$ batched $\mathsf{B2A}$ conversions 
to locally compute the arithmetic mask $\ashare{r}$.
% \begin{equation}
%   \ashare{r}_b:=\sum_{i=0}^{\ell-1}2^i\ashare{r_i}_b\pmod{2^\ell}.
%   \label{eq:drelu-mask}
% \end{equation}
Crucially, these boolean mask bits directly provide $\bshare{r_H}$ and 
the segment shares of $r_L$ (Lines 1--2), enabling the parties 
to set $\bshare{q}_b:=\bshare{r_H}_b\oplus[b{=}0]$ and 
generate the segment keys for the threshold $r_L$ (Line 3--4).

In the online phase, each $P_b$ masks its input to open $\hat{x}=x+r$ in one round and the
segment keys are evaluated at $\hat{y}=\hat{x}\bmod 2^{\ell-1}$ to obtain
$\bshare{\mathit{res}}$. Each party then computes 
the boolean share $\bshare{w}_b$ locally.

\begin{algorithm}[h]
  \caption{Dealerless DReLU ($\Pi_{\DReLU}$)}
  \label{alg:drelu}
  \textbf{Input:} $P_b$ holds the additive share $\ashare{x}_b$ of $x \in \Z_{2^\ell}$. \\
  \textbf{Output:} $P_b$ obtains the boolean share $\bshare{w}_b$ of $w = \DReLU(x)$.

  \vspace{0.3em}
  \textbf{Offline Phase:}
  \begin{algorithmic}[1]
    \State $P_b$ samples $\rho_{i,b}\stackrel{\$}{\leftarrow}\{0,1\}$ for $i=0,\ldots,\ell-1$ and sets $\bshare{r_i}_b:=\rho_{i,b}$.
    \State The parties batch $\ashare{r_i}\gets\mathsf{B2A}(\bshare{r_i})$ and set $\ashare{r}_b:=\sum_{i=0}^{\ell-1}2^i\ashare{r_i}_b\bmod2^\ell$.
    \State $P_b$ sets $\bshare{q}_b:=\bshare{r_{\ell-1}}_b\oplus[b{=}0]$ and takes $\bshare{r_L}_b$ as the lower $\ell-1$ mask bits.
    \State $P_b$ generates the segment keys for threshold $r_L$ via the tree-branching protocols (Section~\ref{sec:core}).
  \end{algorithmic}

  \vspace{0.3em}
  \textbf{Online Phase:}
  \begin{algorithmic}[1]
  \State $P_b$ computes $\ashare{\hat{x}}_b := \ashare{x}_b + \ashare{r}_b$ and opens $\hat{x}$ via $\mathcal{F}_{\mathrm{open}}$.
  \State $P_b$ evaluates the segment keys at $\hat{x}\bmod 2^{\ell-1}$ to obtain $\bshare{\mathit{res}}_b=\bshare{\ind{\hat{y}<r_L}}_b$. \Comment{Grotto-style prefix-parity walk}
  \State $P_b$ computes $u_b := \big( \bshare{q}_b + \bshare{\mathit{res}}_b + [b{=}0]\cdot\mathrm{MSB}(\hat{x}) \big) \bmod 2$. 
  \State $P_b$ outputs $\bshare{w}_b := u_b$.
  \end{algorithmic}
\end{algorithm}

\textbf{Complexity.} Let $m=\ell-1$, let $d=\max_j\wseg_j$ for the
segmentation of $r_L$, and let $C_{\mathrm{key}}(m,\boldsymbol{\wseg})$ be
defined by Equation~\eqref{eq:key-cost}. Preprocessing a Boolean DReLU key
draws on three components that all start from the same shared mask bits and
therefore run concurrently: the $\ell$ parallel
$\mathsf{B2A}:\F_2\rightarrow\Z_{2^\ell}$ calls that assemble $r$ contribute
$\ell(\lambda+\ell)$ bits in one round, segment-key generation for threshold
$r_L$ contributes $C_{\mathrm{key}}(m,\boldsymbol{\wseg})$ bits in $d+2$
rounds, and the $2(\nseg-1)$ bit triples consumed by the balanced
cross-segment reduction contribute $4(\nseg-1)(\lambda+1)$ bits. The gate thus
takes $d+2$ offline rounds with
$\ell(\lambda+\ell)+C_{\mathrm{key}}(m,\boldsymbol{\wseg})+4(\nseg-1)(\lambda+1)$
bits to communicate at the offline stage. Online, opening $\hat{x}$ costs
$\ell$ bits in one round, the $\nseg$ Grotto walks and the closing XORs with
$q$ and $\mathrm{MSB}(\hat{x})$ are local, and the balanced reduction
communicates $4(\nseg-1)$ bits in $\lceil\log_2\nseg\rceil$ rounds, so
evaluation needs $1+\lceil\log_2\nseg\rceil$ rounds with $\ell+4(\nseg-1)$
bits of data exchange. We denote these four totals
$C^{\mathrm{off}}_{\DReLU}$, $R^{\mathrm{off}}_{\DReLU}$,
$C^{\mathrm{on}}_{\DReLU}$ and $R^{\mathrm{on}}_{\DReLU}$, and reuse them for
the derived gates of \S\ref{sec:derived}. For
$(\ell,\lambda,\nseg)=(32,128,2)$ with widths $(16,15)$, this is $18$ offline
rounds and $9{,}828$ offline bits per party, followed by $2$ online rounds and
$36$ online bits per party. Key generation performs $147{,}462$ local hashes
per party, and the two online segment walks perform $33$ local hashes.


\subsection{Lookup Table}\label{sec:lut}
Evaluating nonlinear functions remains a primary performance bottleneck in
secure inference. While spline approximations require 
$\mathcal{O}(k)$ expensive secure comparisons 
and introduce approximation errors~\cite{pika}, 
LUT evaluation is exact for the supplied table; when the table stores a
discretized activation, its function-approximation error must be reported
separately.
Algorithm~\ref{alg:LUT} presents our dealerless LUT
protocol for a $q$-bit index $x\in\Z_{2^q}$ and a public table
$\vec T\in(\Z_{2^p})^{2^q}$.

In the offline phase, $P_0$ and $P_1$ generate 
a random mask $\ashare{r}$ and a shared one-hot 
vector $\vec{v}$ such that $\vec{v}[r]=1$. To avoid the 
prohibitive cost of full-domain expansion, 
the parties efficiently construct $\vec{v}$ via 
the tensor product of $\nseg$ short per-segment DPFs, 
as defined in \eqref{eq:dpf-decomp}. During the online phase, 
the parties reconstruct the masked input $\hat{x} = x - r$. Each $P_b$ 
then locally applies a cyclic shift to $\vec{v}$ by $\hat{x}$ positions, 
yielding an aligned vector $\vec{u}$ such that $\vec{u}[x]=1$. 
Finally, the output share $\ashare{\vec{T}[x]}_b$ is obtained via a 
local inner product between $\vec{u}$ and $\vec{T}$.

\begin{algorithm}[h]
  \caption{Dealerless lookup table ($\Pi_{\mathsf{LUT}}$)}
  \label{alg:LUT}
  \textbf{Input:} $P_b$ holds $\ashare{x}_b$ for $x\in\Z_{2^q}$ and a public table $\vec{T}\in(\Z_{2^p})^{2^q}$. \\
  \textbf{Output:} $P_b$ holds $\ashare{w}_b$ over $\Z_{2^p}$, where $w = \vec{T}[x]$.

  \vspace{0.3em}
  \textbf{[Offline]} $P_0$ and $P_1$ jointly, with no dealer:
  \begin{algorithmic}[1]
    \State $P_b$ samples Boolean mask-bit shares $\bshare{r_i}_b$ for
    $i=0,\ldots,q-1$, batches
    $\ashare{r_i}\gets\mathsf{B2A}(\bshare{r_i})$, and sets
    $\ashare r_b:=\sum_i2^i\ashare{r_i}_b\bmod 2^q$.
    \State Run the dealerless DPF generation of Section \S\ref{sec:core} for a 
    payload of $1$ at the secret index $r$, obtaining the $\nseg$ per-segment DPF key shares $k_b$.
    \State Upon receiving $k_b$, the parties perform a batched
    evaluation across the segment-domain to obtain the shared 
    vector $\ashare{\vec{v}}$, where 
    $\vec{v}[i] \leftarrow \Pi_{\mathrm{Eval}}(k_b, i)$ 
    for $i \in \{0, \dots, 2^q-1\}$.
  \end{algorithmic}

  \vspace{0.3em}
  \textbf{[Online]} on input $\ashare{x}$, each $P_b$:
  \begin{algorithmic}[1]
    \State $\ashare{\hat{x}}_b \gets \ashare{x}_b - \ashare{r}_b$, then open $\hat{x}$ via $\mathcal{F}_{\mathrm{open}}$. \Comment{1 round}
    \State $\ashare{\vec u}_b \gets \ashare{\vec v}_b \ggg \hat{x}$ \quad and \quad $\ashare{w}_b \gets \sum_{i=0}^{2^q-1} \vec{T}[i]\cdot\ashare{\vec u[i]}_b$. \Comment{local shift and inner product}
    \State \textbf{return} $\ashare{w}_b$
  \end{algorithmic}
\end{algorithm}

\textbf{Complexity.} Let the LUT index have $q$ bits, the table entries and
output have $p$ bits, and let
$d=\max_j\wseg_j$. Generating its segment keys costs
$R_{\mathrm{key}}=d+2$ offline rounds,
$C_{\mathrm{key}}(q,\boldsymbol{\wseg})$ bits per party, and
$\sum_j(3\cdot2^{\wseg_j-1}+3)$ hashes. Expanding the full segment vectors
uses $\sum_j(3\cdot2^{\wseg_j-1}-1)$ further local hashes. Fix a balanced binary tensor tree
$\mathcal T$ over the segments. For an internal node $v$, let
$W_v=\sum_{j\in\operatorname{leaves}(v)}\wseg_j$ and define
$A_{\mathcal T}=\sum_{v\in\operatorname{int}(\mathcal T)}2^{W_v}$.
Materializing the Boolean one-hot vector requires exactly $A_{\mathcal T}$
secure \textsc{and} gates. It therefore costs
$2A_{\mathcal T}(\lambda+2)$ offline bits, including triple generation and
consumption, in $\lceil\log_2\nseg\rceil$ dependent rounds after key
generation. As the result is used in an arithmetic inner product, converting
its $2^q$ entries adds $2^q(\lambda+p)$ bits and one round. The $q$ mask bits
also require $q(\lambda+q)$ bits, batched with key generation. Hence
\begin{equation}
\begin{aligned}
C^{\mathrm{off}}_{\mathsf{LUT}}(q,p)
 &=q(\lambda+q)+C_{\mathrm{key}}(q,\boldsymbol{\wseg})\\
 &\quad+2A_{\mathcal T}(\lambda+2)+2^q(\lambda+p),\\
R^{\mathrm{off}}_{\mathsf{LUT}}(q,p)&=d+3+\lceil\log_2\nseg\rceil,
\end{aligned}
\label{eq:lut-offline-cost}
\end{equation}
where all independent preprocessing calls are batched. Online, opening the
masked index costs $q$ bits in one round; the rotation and public-table inner
product are local:
\begin{equation}
C^{\mathrm{on}}_{\mathsf{LUT}}(q,p)=q,\qquad
R^{\mathrm{on}}_{\mathsf{LUT}}(q,p)=1.
\label{eq:lut-online-cost}
\end{equation}
This accounting makes the scope explicit: segmentation reduces dealerless key
generation, but materializing a full $2^q$ one-hot vector remains exponential
in the reduced LUT index width $q$.

\subsection{Dealerless Truncation}\label{sec:trunc}
Every fixed-point product leaves $2f$ fractional bits, so $\Fmul$ is followed
by a rescaling by $f$ bits. We use SIGMA's exact truncate-reduce
protocol~\cite{sigma}. In preprocessing, the parties sample a uniform mask
$r=2^f h+l$ in $\mathbb Z_{2^\ell}$, where $0\le l<2^f$. They hold arithmetic
shares of $h$ and $r$, and Boolean shares of the bits of $l$. A batch of
$f$ Boolean-to-arithmetic conversions constructs the arithmetic shares of $r$.
The parties generate a one-bit correlated-DPF comparison key for the secret
threshold $l$. They also prepare Boolean and arithmetic shares of an independent
random bit $c$.

Online, the parties open $\hat z=z+r\bmod 2^\ell$. The DPF comparison returns
Boolean shares $w_b$ of $w=\ind{(\hat z\bmod 2^f)<l}$. They open only the
masked bit $\bar w=\bigoplus_b(w_b\oplus c_b)$ and use the arithmetic shares
$C_b$ of $c$ to recover an arithmetic share of $w$:
\begin{equation}
\begin{aligned}
u_b&=b\bar w+C_b-2\bar w C_b,\\
\ashare{z'}_b
 &=b\left\lfloor\hat z/2^f\right\rfloor-h_b-u_b
   \pmod {2^{\ell-f}}.
\end{aligned}
\label{eq:trunc}
\end{equation}
Here $b\in\{0,1\}$ indexes the parties and $h_b$ is their share of $h$.
The low-bit comparison removes the borrow caused by subtracting $r$ from
$\hat z$. Thus $z'=\lfloor z/2^f\rfloor\bmod 2^{\ell-f}$ for every input,
including negative two's-complement values. We also use this truncation to
form the bounded SIGMA LUT index.

\textbf{Complexity.} Let $C^{\mathrm{off}}_{\mathsf{B2A}}(f+1,\ell)$ denote
the cost of converting the $f$ low mask bits and $c$ in the $\ell$-bit ring.
Let $C^{\mathrm{off}}_{\mathsf{DPF}}(f,1)$ denote dealerless generation of
the one-bit DPF comparison key. The analogous $R$ terms count offline rounds.
Our implementation runs these steps in sequence. Per party, the logical
communication costs are
\begin{equation}
\begin{aligned}
C^{\mathrm{off}}_{\mathsf{Trunc}}(f)
 &=C^{\mathrm{off}}_{\mathsf{B2A}}(f+1,\ell)
   +C^{\mathrm{off}}_{\mathsf{DPF}}(f,1),\\
R^{\mathrm{off}}_{\mathsf{Trunc}}(f)
 &=R^{\mathrm{off}}_{\mathsf{B2A}}(f+1,\ell)
   +R^{\mathrm{off}}_{\mathsf{DPF}}(f,1),\\
C^{\mathrm{on}}_{\mathsf{Trunc}}(f)&=\ell+1,\qquad
R^{\mathrm{on}}_{\mathsf{Trunc}}(f)=2.
\end{aligned}
\label{eq:trunc-cost}
\end{equation}
SIGMA counts one round and two bits in total for the comparison when $\hat z$
is already public. Our arithmetic-share interface first opens $\hat z$ and
then opens the masked comparison bit. In measured communication, the latter
occupies one byte per party; the ring opening uses one, two, four, or eight
bytes for $\ell\leq8,16,32,64$, respectively.

We use $\MulTr(f)=\Trunc(f)\circ\Fmul$ below for a fixed-point
multiplication followed by exact $f$-bit rescaling. Combining the Beaver cost
from \S\ref{sec:prelim} with~\eqref{eq:trunc-cost} gives the following costs.
Here $R^{\mathrm{off}}_{\mathsf{Mul}}(\ell)$ counts the rounds used to generate
one dealerless Beaver triple:
\begin{equation}
\begin{aligned}
C^{\mathrm{off}}_{\mathsf{MT}}(f)
 &=2\ell(\lambda+\ell)+C^{\mathrm{off}}_{\mathsf{Trunc}}(f),\\
R^{\mathrm{off}}_{\mathsf{MT}}(f)
 &=R^{\mathrm{off}}_{\mathsf{Mul}}(\ell)
   +R^{\mathrm{off}}_{\mathsf{Trunc}}(f),\\
C^{\mathrm{on}}_{\mathsf{MT}}(f)
 &=3\ell+1,\\
R^{\mathrm{on}}_{\mathsf{MT}}(f)
 &=3.
\end{aligned}
\label{eq:multr-cost}
\end{equation}
The Beaver opening precedes the input-mask opening and DPF comparison.

\subsection{Derived Nonlinear Gates}\label{sec:derived}
Based on $\Pi_{\text{DReLU}}$, we first construct Max and ReLU functions.
We then build inversion, reciprocal square root, GELU, Sigmoid, Tanh, Softmax,
and LayerNorm from $\Pi_{\text{LUT}}$ with $\Pi_{\text{DReLU}}$.
These protocols complete the nonlinear operations used by neural-network inference~\cite{cryptflow2,sigma,falcon}
and are detailed in Appendix \ref{appendix:Secure 3PC Protocols}. 

\paragraph*{ReLU}
Given a shared value $x$ and a target value $res$, we first extend $\Pi_{\text{DReLU}}$ to 
a secure selection protocol defined as $\text{Select}(x, res) = 1\{x \geq 0\} \cdot res$ ($\Pi_{\text{Select}}$). 
The boolean
DReLU result serves as a direct input to the mixed-share multiplication.
Consequently, the ReLU function can be defined as $\text{ReLU}(x) = \text{Select}(x, x)$.
Writing $C^{\mathrm{off}}_{\DReLU}$ and $C^{\mathrm{on}}_{\DReLU}$ for the two
DReLU totals of \S\ref{sec:Drelu}, one $\ell$-bit ReLU costs
\begin{equation}
\begin{aligned}
C^{\mathrm{off}}_{\ReLU}&=C^{\mathrm{off}}_{\DReLU}
 +4\ell(\lambda+\ell+1),\\
R^{\mathrm{off}}_{\ReLU}&=\max\{R^{\mathrm{off}}_{\DReLU},3\},\\
C^{\mathrm{on}}_{\ReLU}&=C^{\mathrm{on}}_{\DReLU}+\ell+1,
\\R^{\mathrm{on}}_{\ReLU}&=2+\lceil\log_2\nseg\rceil.
\end{aligned}
\label{eq:select-cost}
\end{equation}
The added terms are exactly one $\BitBA$ preprocessing and one masked
ring-and-bit opening. An $\ell$-bit $\Pi_{\mathsf{Select}}$ has the same four
costs as~\eqref{eq:select-cost}; replacing the selected value $x$ by another
$\ell$-bit arithmetic share does not change the protocol.

\paragraph*{MaxPool and ArgMax}
To evaluate $k$ shared values, we structure
pairwise comparisons as a balanced binary tree.
For candidates $(a, i_a)$ and $(b, i_b)$, we define
$\text{Max}(a, b) = \text{Select}(a - b, a - b) + b,
\text{ArgMax}(i_a, i_b) = \text{Select}(a - b, i_a - i_b) + i_b$.
For MaxPool, every internal node has the ReLU cost above. A binary tree over
$k$ inputs has $k-1$ internal nodes, and calls at one level are batched. Thus
\begin{equation}
\begin{aligned}
C^{\mathrm{off}}_{\mathsf{MaxPool}}(k)&=(k-1)C^{\mathrm{off}}_{\ReLU},\\
R^{\mathrm{off}}_{\mathsf{MaxPool}}(k)&=R^{\mathrm{off}}_{\ReLU},\\
C^{\mathrm{on}}_{\mathsf{MaxPool}}(k)&=(k-1)C^{\mathrm{on}}_{\ReLU},\\
R^{\mathrm{on}}_{\mathsf{MaxPool}}(k)
 &=\lceil\log_2 k\rceil R^{\mathrm{on}}_{\ReLU}.
\end{aligned}
\label{eq:maxpool-cost}
\end{equation}
ArgMax reuses each comparison but selects both the value and its index. For a
$p$-bit index, the second selection at every internal node adds one $\BitBA$.
The value and index selections run in parallel, so
\begin{equation}
\begin{aligned}
C^{\mathrm{off}}_{\mathsf{ArgMax}}(k,p)
 &=(k-1)\bigl(C^{\mathrm{off}}_{\ReLU}
       +4p(\lambda+p+1)\bigr),\\
R^{\mathrm{off}}_{\mathsf{ArgMax}}(k,p)
 &=\max\{R^{\mathrm{off}}_{\ReLU},3\},\\
C^{\mathrm{on}}_{\mathsf{ArgMax}}(k,p)
 &=(k-1)\bigl(C^{\mathrm{on}}_{\ReLU}+p+1\bigr),\\
R^{\mathrm{on}}_{\mathsf{ArgMax}}(k,p)
 &=\lceil\log_2 k\rceil R^{\mathrm{on}}_{\ReLU}.
\end{aligned}
\label{eq:argmax-cost}
\end{equation}

\paragraph*{Sigmoid and Tanh}
Recurrent gates need $\sigma$ and $\tanh$ rather than $\ReLU$. For sigmoid,
we use the folded approximation in Appendix~G of SIGMA~\cite{sigma}, which
exploits the even, rapidly decaying correction
\begin{equation}
  \delta(x)=\frac{1}{1+e^{|x|}}, \qquad
  \sigma(x)=\delta(x)+\ind{x\geq 0}\bigl(1-2\delta(x)\bigr).
  \label{eq:sigmoid-decomp}
\end{equation}
We instantiate this approximation with $f=12$ fractional bits and a $q=10$
bit table index. Let $x_f=\lfloor 2^f x\rfloor$ be the signed ring encoding
and set $B_0=2^{f+4}$ and $B_f=B_0-1$. For $f=12$, the correction is
below one output ULP once $|x|\geq 16$, so clipping at $B_f$ preserves the
encoded endpoint. The secure evaluation first obtains
$d=\DReLU(x_f)=\ind{x_f\geq 0}$, computes
\begin{equation}
\begin{aligned}
  p&=\BitBA(d,x_f), & a&=2p-x_f,\\
  i&=\ind{a<B_0}, & c&=B_f+i(a-B_f).
\end{aligned}
\label{eq:sigmoid-fold}
\end{equation}
where $a=|x_f|$ and $c=\min(a,B_f)$. The clipping comparison is
$i=\DReLU(a-B_0)\oplus 1$. This order computes the magnitude and clips it
before truncating the LUT index. Since $c$ is
nonnegative, truncation by $f-6$ bits produces the ten-bit index
\begin{equation}
  t=\left\lfloor\frac{c}{2^{f-6}}\right\rfloor, \qquad
  T_\delta[j]=\left\lfloor
    \frac{2^f}{1+\exp(j/2^6)}\right\rfloor,
  \quad 0\leq j<2^{10}.
  \label{eq:sigmoid-table}
\end{equation}
The table lookup returns $\widehat\delta=T_\delta[t]$. The encoded sigmoid is
then reconstructed with one mixed-share product,
\begin{equation}
  \widehat\sigma_f(x)=\widehat\delta
    +\BitBA\bigl(d,2^f-2\widehat\delta\bigr).
  \label{eq:sigmoid-stitch}
\end{equation}
Thus the table contains $2^{10}=1024$ entries with grid spacing $2^{-6}$,
independent of the 32-bit ring width. At $f=12$, inputs at or beyond
$\lvert x\rvert=16$ map to $\widehat\delta=0$ and therefore to the appropriate
sigmoid endpoint. The lookup is exact for this integer table; the resulting
activation is an approximation of the real sigmoid.

Our C++ implementation realizes the two sign tests with the
$\textsf{BitDec}$ primitive, the three mixed products with
$\textsf{BitXA}$ (denoted $\BitBA$ here), the index truncation with
$\Pi_{\mathsf{Trunc}}$, and the final table evaluation with
$\Pi_{\mathsf{LUT}}$. Consequently, one scalar activation uses one public
$q=10$ LUT key, one six-bit truncation record, and three $\BitBA$ preprocessing
records; online it performs two $\textsf{BitDec}$ calls, three mixed products,
one truncation, and one LUT evaluation. The corresponding communication terms are
\begin{equation}
\begin{aligned}
C^{\mathrm{off}}_{\mathsf{Sigmoid}}
 &= C^{\mathrm{off}}_{\mathsf{LUT}}(10,\ell)
  + C^{\mathrm{off}}_{\mathsf{Trunc}}(6)
  + 3C^{\mathrm{off}}_{\BitBA},\\
C^{\mathrm{on}}_{\mathsf{Sigmoid}}
 &= 2C^{\mathrm{on}}_{\textsf{BitDec}}(\ell)
  + 3C^{\mathrm{on}}_{\BitBA}(\ell)\\
 &\quad+ C^{\mathrm{on}}_{\mathsf{Trunc}}(6)
  + C^{\mathrm{on}}_{\mathsf{LUT}}(10,\ell).
\end{aligned}
\label{eq:sigmoid-cost}
\end{equation}
In the current implementation, the Boolean triples consumed inside
$\textsf{BitDec}$ are generated by its evaluation calls and are included in
$C^{\mathrm{on}}_{\textsf{BitDec}}$; the three $\BitBA$ records, truncation record,
and LUT key are the explicit preprocessing material above.
This activation uses a common fractional scale at its input and output, and its
sign tests do not call the segmented comparison path. It is a separate numerical
construction from SiRNN's sigmoid/tanh approximations; replacing those functions
with it would change the integer program.
The vector interface invokes this scalar protocol independently for each
element. Within one activation, the sign, clipping, lookup, and branch-selection
dependencies form the online chain used for its round count.

We derive a tanh approximation from this sigmoid using the identity
$\tanh(x)=2\sigma(2x)-1$, reusing the same table:
\begin{equation}
  \widehat{\tanh}_f(x)=2\widehat\sigma_f(2x)-2^f.
  \label{eq:tanh-cost}
\end{equation}
This tanh construction extends SIGMA's sigmoid approximation.
The input doubling and the subtraction are local shifts and public constants,
so $\tanh$ has the same communication terms as one sigmoid. The doubled input
reaches the clipping bound at $|x|=8$, which is the effective unsaturated range
of this construction. The implementation requires $\ell>f+4$ so that the
encoded $B=2^{f+4}$ bound is representable; for $f=12$ this means
$\ell\geq17$. Because $2x$ is formed in the input ring, the tanh preactivation
must also stay within the ring headroom required by this doubling.

\begin{table}[t]
\centering
\caption{Measured error of the SIGMA-based sigmoid and the derived tanh at
$q=10$ and
$f=12$. The input scan covers every representable value in $[-16,16)$;
ULP is $2^{-12}$ at the output.}
\label{tab:sigma-error}
\resizebox{\columnwidth}{!}{%
\begin{tabular}{lrrr}
\toprule
Function & Max. absolute error & Max. ULP error & Input at maximum\\
\midrule
$\sigma$ & $0.0038451390$ & $15.7497$ & $-0.015380859375$\\
$\tanh$ & $0.0075682149$ & $30.9994$ & $-0.007568359375$\\
\bottomrule
\end{tabular}
}
\end{table}
The errors in Table~\ref{tab:sigma-error} compare the integer LUT output with
the mathematical functions, not with an ideal public lookup. They were obtained
by exhaustive enumeration in the precision mirror. The two-party correctness
test checks representative negative, positive, and saturation points through the
same secure evaluation path.

\paragraph*{Recurrent cells}
At the protocol level, these gates compose into a recurrent step. An LSTM
cell evaluates its four pre-activations as one batched affine map, applies
$\Pi_{\text{Sigmoid}}$ to three of them and $\tanh$ to the fourth in a single
batch, and then updates state through
$c_t=f_t\odot c_{t-1}+i_t\odot g_t$ and $h_t=o_t\odot\tanh(c_t)$. The
Hadamard products are secret-by-secret, so each is one $\Fmul$ followed by one
$\Trunc$. This is why the recurrent case needs the arithmetic gates of
\S\ref{sec:trunc} and not only the comparison gate that suffices for $\ReLU$
networks. GRU cells differ only in which products appear. Batching within a step
makes its round count independent of the hidden width; across steps the
recurrence is inherently sequential, so a $T$-step sequence costs $T$ times the
per-step depth. State is the only value that persists between steps, and it
persists as shares, so no intermediate is ever opened.

For a hidden vector of width $h$, count the formulas below per coordinate.
The four pre-activation gates contribute four parallel Sigmoid-equivalent
calls, the state update uses two parallel $\MulTr$ calls, and computing
$h_t$ adds one Tanh-equivalent call followed by one $\MulTr$. Therefore one
LSTM step communicates
\begin{equation}
\begin{aligned}
C^{\mathrm{off}}_{\mathsf{LSTM}}(h)
 &=h\bigl(5C^{\mathrm{off}}_{\mathsf{Sigmoid}}
          +3C^{\mathrm{off}}_{\mathsf{MT}}(f)\bigr),\\
R^{\mathrm{off}}_{\mathsf{LSTM}}
 &=\max\{R^{\mathrm{off}}_{\mathsf{Sigmoid}},
          R^{\mathrm{off}}_{\mathsf{MT}}(f)\},\\
C^{\mathrm{on}}_{\mathsf{LSTM}}(h)
 &=h\bigl(5C^{\mathrm{on}}_{\mathsf{Sigmoid}}
          +3C^{\mathrm{on}}_{\mathsf{MT}}(f)\bigr),\\
R^{\mathrm{on}}_{\mathsf{LSTM}}
 &=2R^{\mathrm{on}}_{\mathsf{Sigmoid}}
   +2R^{\mathrm{on}}_{\mathsf{MT}}(f).
\end{aligned}
\label{eq:lstm-cost}
\end{equation}
The width multiplies communication but not rounds. A sequence of $T$ steps
multiplies both communication totals and the online depth in
\eqref{eq:lstm-cost} by $T$; preprocessing for all steps remains batchable.

\textbf{Inverse and Reciprocal Square Root.} 
Following the domain reduction strategy from Mosformer~\cite{Mosformer}, 
our protocol evaluates $\Pi_{\mathsf{Inv}}$ and $\Pi_{\mathsf{Rsqrt}}$
using $\Pi_{\text{LUT}}$. Specifically, an interval 
selection protocol identifies the distribution range of $x$, 
mapping it linearly to a compact subdomain $[1/b, 1]$. The 
resulting lookup table requires $(1 - 1/b) \cdot 2^f$ entries, where $f$
denotes the target precision. Equations~\eqref{eq:inverse-cost} and
\eqref{eq:rsqrt-cost} give the exact costs.

% \textbf{GELU.} 
% The GELU function is defined as 
% $\text{GELU}(x) = 0.5x \left(1 + \tanh\left(\sqrt{\frac{2}{\pi}}(x + 0.044715x^3)\right)\right)$. 
% To securely evaluate $\tanh$, we adapt 
% the approach of SIGMA~\cite{sigma} by implementing a range-limited lookup table via $\Pi_{\text{LUT}}$.


\textbf{Exponentiation.} We instantiate the exponentiation construction of
BOLT~\cite{Bolt} with $\Pi_{\mathsf{Select}}$ and $\Pi_{\mathsf{LUT}}$.
For the one-Select/one-LUT instantiation used here, let $q_e$ be the reduced
LUT index width. The selection determines the public-range branch and the LUT
then evaluates that branch, so
\begin{equation}
\begin{aligned}
C^{\mathrm{off}}_{\mathsf{Exp}}(q_e)
 &=C^{\mathrm{off}}_{\ReLU}+C^{\mathrm{off}}_{\mathsf{LUT}}(q_e,\ell),\\
R^{\mathrm{off}}_{\mathsf{Exp}}(q_e)
 &=\max\{R^{\mathrm{off}}_{\ReLU},R^{\mathrm{off}}_{\mathsf{LUT}}(q_e,\ell)\},\\
C^{\mathrm{on}}_{\mathsf{Exp}}(q_e)
 &=C^{\mathrm{on}}_{\ReLU}+q_e,\\
R^{\mathrm{on}}_{\mathsf{Exp}}(q_e)
 &=R^{\mathrm{on}}_{\ReLU}+1.
\end{aligned}
\label{eq:exp-cost}
\end{equation}
The offline calls are independent and batch, whereas the online LUT consumes
the selected index and therefore follows the selection.

% For an input vector $\vec{x} \in \mathbb{R}^k$, 
% the function is defined as 
% $\text{softmax}(\vec{x})[i] = \frac{\exp(x[i] - x^{\text{max}})}{\sum_{j=0}^{k-1} \exp(x[j] - x^{\text{max}})}$, 
% where $x^{\text{max}}$ is the maximum element in $\vec{x}$. 
% Our secure softmax computation proceeds in three phases: 
% (i) \emph{Maximum selection}: Computing $x^{\text{max}}$ via $\lceil \log_2 k \rceil$ 
% interactive rounds using $\Pi_{\text{Select}}$; (ii) \emph{Exponentiation}: 
% Evaluating the exponential terms using BOLT's approximation~\cite{Bolt}, combining 
% $\Pi_{\text{Select}}$ and $\Pi_{\text{LUT}}$; and (iii) \emph{Division}.

% Given an input vector $\vec{x} \in R^k$, 
% the function is defined as: 
% $\text{softmax}(x)[i] = \frac{\exp\left(x[i] - x^{max} \right)}{\sum_{j=0}^{k-1} \exp\left(x[j] - x^{max} \right)}$, 
% where $x^{max}$ is the max element of the input vector $\vec{x}$. 
% As the main component of the self-attention mechanism, 
% the secure softmax computation consists of three modules: 
% $(\text{i})$ the maximum value selection module. 
% We calculate the maximum through $\lceil log_2^k \rceil$ rounds of 
% interaction using the $\Pi_{\text{Select}}$. 
% $(\text{ii})$ the exponentiation module. This module 
% adopts the approximation method proposed by BOLT \cite{Bolt}, 
% implemented collaboratively through $\Pi_{\text{Select}}$ and $\Pi_{\text{LUT}}$. 
% $(\text{iii})$ the division module. 

\subsection{The Protocol \texorpdfstring{$\Pi_{\mathsf{Infer}}$}{Pi-Infer}}\label{sec:infer-proto}
Algorithm~\ref{alg:infer} describes the comparison gate used in an inference
composition. Private linear layers require secure multiplication and rescaling
in addition to local arithmetic. Each comparison gate runs the blind and open step, the parallel
segment evaluation, and the multiplication-scan stitch. Reductions such as maxpool
and argmax fold pairwise comparisons, and a table or spline activation runs
$\nseg$ parallel reduced-domain lookups. All comparison gates of a layer are
batched, so a layer costs a constant number of rounds regardless of width.

\begin{algorithm}[t]
\caption{One comparison gate ($\Pi_{\mathsf{Infer}}$)}
\label{alg:infer}
\textbf{Input:} $P_b$ holds $\ashare{x}_b$, the preprocessed segment keys $\{k\seg{j}\}_j$, and $\ashare{r_{\mathrm{in}}}_b$. \\
\textbf{Output:} $P_b$ holds $\bshare{z}_b=\bshare{\ind{x<c}}_b$, for example the $\DReLU$ bit.

\vspace{0.3em}
\textbf{[Online]} on input $\ashare{x}$, each $P_b$:
\begin{algorithmic}[1]
\State $\ashare{\xhat}_b\gets\ashare{x}_b+\ashare{r_{\mathrm{in}}}_b$, then open $\xhat$ \Comment{1 round, $\ell$ bits}
\State locally split $\xhat\to(\xhat\seg{1},\dots,\xhat\seg{\nseg})$
\State \textbf{in parallel for} $j\in[\nseg]$: $\bshare{L_j}(,\bshare{e_j})\gets\Eval(b,k\seg{j},\xhat\seg{j})$ \Comment{Grotto-style prefix-parity walk}
\State prefix-\textsc{and} scan $\bshare{E_{j+1}}\gets\bshare{E_j}\wedge\bshare{e_j}$, then $\bshare{z}\gets\bigoplus_j \bshare{E_j}\wedge\bshare{L_j}$ \Comment{$\bigO(\log\nseg)$ rounds}
\State apply wrap-correction (local XOR), then $\mathsf{B2A}$ if an arithmetic result is needed
\State \textbf{return} $\bshare{z}_b=\bshare{\ind{x<c}}_b$
\end{algorithmic}
\end{algorithm}

\subsection{Theoretical Analysis}
\paragraph*{Online cost}
For a Boolean comparison on an $\ell$-bit input, opening $\xhat$ costs $\ell$
bits in one round, the $\nseg$ Grotto walks communicate nothing, and the
$2(\nseg-1)$ scan \textsc{and}s cost $4(\nseg-1)$ bits in
$\lceil\log_2\nseg\rceil$ rounds. The exact total is therefore
\begin{equation}
C^{\mathrm{on}}_{\mathsf{Cmp}}=\ell+4(\nseg-1),\qquad
R^{\mathrm{on}}_{\mathsf{Cmp}}=1+\lceil\log_2\nseg\rceil.
\end{equation}
An arithmetic comparison output adds one
$\mathsf{B2A}:\F_2\rightarrow\Z_{2^p}$, namely $\lambda+p$ bits and one round.
Its exact online total is therefore
\begin{equation}
C^{\mathrm{on}}_{\mathsf{CmpA}}=\ell+4(\nseg-1)+\lambda+p,qquad
R^{\mathrm{on}}_{\mathsf{CmpA}}=2+\lceil\log_2\nseg\rceil.
\end{equation}
For a batch of $B$ independent comparisons, communication becomes
$B C^{\mathrm{on}}_{\mathsf{Cmp}}$, whereas the round count is unchanged. An
$\ell$-bit $\Fmul$ similarly costs $2\ell$ online bits in one round; the exact
truncation and activation costs are stated in their respective subsections.

\paragraph*{Offline cost}
When Boolean shares of the threshold bits are already available, one comparison
key costs $C_{\mathrm{key}}(\ell,\boldsymbol{\wseg})$ bits and $d+2$ rounds,
where $d=\max_j\wseg_j$. Its scan triples add
$4(\nseg-1)(\lambda+1)$ bits but no round because their COT generation is
batched with $\PREP$. Thus
\begin{equation}
C^{\mathrm{off}}_{\mathsf{Cmp}}
=C_{\mathrm{key}}(\ell,\boldsymbol{\wseg})
 +4(\nseg-1)(\lambda+1),\qquad
R^{\mathrm{off}}_{\mathsf{Cmp}}=d+2.
\end{equation}
If the threshold is initially arithmetic-shared, $\Pi_{\mathsf{Split}}$ adds
$(2\ell-3)(2\lambda+4)$ bits. The first random-COT round of $\PREP$ overlaps
the triple-generation round of Split, but retargeting and the $d$
correction-word rounds wait for the Boolean threshold bits, so the two depths
add rather than overlap:
\begin{equation}
C^{\mathrm{off}}_{\mathsf{CmpAThr}}
=C^{\mathrm{off}}_{\mathsf{Cmp}}+(2\ell-3)(2\lambda+4),\qquad
R^{\mathrm{off}}_{\mathsf{CmpAThr}}=\ell+d+1.
\end{equation}
A batch of $B$ gates multiplies the communication and hash counts by $B$ while
preserving these round counts. The cGGM work is
$\sum_j(3\cdot2^{\wseg_j-1}+3)$ hashes per gate; for uniform widths this is
$\nseg(3\cdot2^{\ell/\nseg-1}+3)$.

For a complete inference circuit, let $B_{t,g}$ be the number of independent
instances of primitive $g$ in nonlinear layer $t$, using a nonoverlapping gate
decomposition. Using the costs above,
\begin{equation}
\begin{aligned}
C^{\mathrm{off}}_{\mathsf{Infer}}
 &=\sum_t\sum_g B_{t,g}C^{\mathrm{off}}_g,&
R^{\mathrm{off}}_{\mathsf{Infer}}
 &=\max_{t,g:B_{t,g}>0}R^{\mathrm{off}}_g,\\
C^{\mathrm{on}}_{\mathsf{Infer}}
 &=\sum_t\sum_g B_{t,g}C^{\mathrm{on}}_g,&
R^{\mathrm{on}}_{\mathsf{Infer}}
 &=\sum_t\max_{g:B_{t,g}>0}R^{\mathrm{on}}_g.
\end{aligned}
\label{eq:infer-cost}
\end{equation}
All gates in one layer batch, while data dependencies between layers add their
online depths. Input-independent preprocessing for the entire circuit batches
globally, so only its maximum round depth remains.

\paragraph*{Security}
\begin{proposition}
$\Pi_{\mathsf{Infer}}$ securely realizes secure comparison/$\DReLU$ against a semi-honest adversary corrupting one party, in the $(\Fcot,\Frand)$-hybrid
model.
\end{proposition}
\begin{IEEEproof}[Proof sketch]
The only opened value is $\xhat=x+r_{\mathrm{in}}$, a one-time pad of $x$, which the simulator samples uniformly, and the public digits $\xhat\seg{j}$ are a
deterministic local function of it and add no leakage. The $\PREP$ transcript is
simulated from $\Fcot$, $\Frand$ outputs are public coins, per-level correction
words are reconstructed from shares uniform given the view of one party, and segment
$\Eval$ is provenance-agnostic prefix-parity on a key the simulator already produces.
The multiplication scan, $\CCMP$, $\mathsf{B2A}$, and $\BitBA$ are simulated
from their standard functionalities; the last uses the one-round construction
reviewed in \S\ref{sec:prelim}. Composing the simulators yields a view indistinguishable from
the real execution.
\end{IEEEproof}

% =====================================================================
\section{Evaluation}\label{sec:eval}

The available experiments characterize the cost of segmented comparison
primitives, the numerical effects of fixed-point formats, the comparison domains
used in mixed-precision rescaling, and the consistency of a historical Google-30
SCI replay. The primitive timings come from two-party C++ executions on localhost.
The ImageNette accuracy results
come from cleartext integer inference, including a CUDA wide-accumulator path.
The historical SCI replay executes 30 complete FastGRNN sequences securely.
These measurements do not yet provide secure end-to-end costs for SCI + ours.

The local CNN/FastGRNN/LSTM timing rows marked \texttt{structure\_only} execute
a cleartext fixed-point runtime. Rows marked \texttt{smoke} check program
execution. Neither supplies secure offline/online latency or communication;
those fields remain unavailable until the corresponding private executor runs.
We first present numerical evidence and primitive measurements, then specify
the end-to-end comparison that remains to be carried out.

\paragraph*{Fixed-point numerical sensitivity}
We distinguish four numerical models of fixed-point inference. E0-B encodes
values as $\lfloor 2^s z\rfloor\bmod 2^\ell$ and reduces stored additions
in $\mathbb Z_{2^\ell}$. Linear and convolution layers accumulate products
in an $(\ell+s)$-bit temporary ring, apply one signed floor truncation by
$s$ after each dot product, and return the output to $\ell$ bits.
Elementwise products are truncated individually. ReLU selects the encoded
input or zero using the ring sign bit. In the reported LSTM precision scan,
sigmoid and tanh are evaluated as real functions and immediately re-encoded;
this scan measures arithmetic quantization with ideal activation gates.
The $q=10$, $f=12$ SIGMA-based activation errors in
Table~\ref{tab:sigma-error} are measured separately.

The strict-ring study, E0-U-R, uses one signed Q$(\ell,s)$ format for
all inputs, weights, biases, inference-time BN affine coefficients, and
operator outputs. Let $W_\ell(z)$ denote reduction modulo $2^\ell$
followed by interpretation in $[-2^{\ell-1},2^{\ell-1}-1]$.
Encoding is $W_\ell(\lfloor 2^s z\rfloor)$. Each scalar product,
accumulator addition, bias lift, and residual addition is reduced in
$\mathbb Z_{2^\ell}$. A convolution or linear output is therefore
\[
q_y=W_\ell\!\left(\left\lfloor\frac{
W_\ell(\sum_i q_{x_i}q_{w_i}+2^s q_b)}{2^s}\right\rfloor\right).
\]
The numerator represents ring additions of ring products; signed floor
rescaling occurs once after the complete dot product. BN affine maps
follow the same rule. Global average pooling first reduces its sum in
the same ring, then floor-divides its signed representative by the public
window size. ReLU, max pooling, and argmax compare signed ring codes.
The logical accumulator width equals $\ell$. A host integer container or
integer limb decomposition implements these modular operations without
retaining out-of-ring high bits for subsequent rescaling.

The historical saturating control, E0-U-S, also uses uniform
Q$(\ell,s)$ storage, but floors and saturates encoding and writeback.
Its convolution and linear layers accumulate unscaled integer products
and lifted biases in overflow-checked 128-bit temporaries, then rescale
once. Residual additions saturate at writeback; global average pooling
divides the wide sum before saturation. E0-U-S separates storage
quantization from accumulator overflow and has different arithmetic from
E0-U-R.

E0-Q-fixed uses a uniform signed Q$(\ell,s)$ format for the entire network.
Every input, quantized weight, bias, BN affine coefficient, intermediate
tensor, and output score is an integer code in
$[-2^{\ell-1},2^{\ell-1}-1]$ with scale $2^{-s}$. Thus $s$ controls the
fractional resolution and $\ell-s$ counts the remaining bits, including
the sign. Each scalar product is formed in a host temporary, floor-rescaled
immediately by $s$, and saturated back to $\ell$ bits. Each dot-product term
is then accumulated with an $\ell$-bit saturating addition. Global average
pooling also saturates each addition before division by the public window
size. BN affine maps, residual additions, pooling, ReLU, logits, and argmax
remain in the integer domain; floating point is used to encode constants
and to evaluate the reference model. This study measures the combined
effects of uniform storage, product rescaling, and finite-width accumulation,
and serves as a per-term saturating-accumulation control.

The uniform-study protocol enumerates all legal
$s\in\{0,\ldots,\ell-1\}$ for each $\ell\in\{12,16,20,24,28,32\}$
on an independent 512-example calibration subset,
selects the highest-accuracy $s$ (ties use the smaller $s$), and evaluates the
selected pair on the complete test split. The original E0-B references are
97.48\% for MLP on MNIST, 88.22\% for ResNet20 on CIFAR-10, and 95.90\% for
the synthetic LSTM task. The uniform studies use ImageNette, a 10-class
ImageNet subset, with 3,925 validation images and ImageNet-1K-pretrained
ResNet-50/101/152 models. They retain the original 1,000-class output heads
and use the same class-balanced training calibration subset. For the E0-U
variants, the measurement protocol includes classification accuracy and
per-example prediction agreement with FP32; validation images do not
participate in scale selection.

The earlier E0-B results on the reference tasks are as follows; parentheses
give the training-selected fractional width.
\begin{center}
\small
\resizebox{\columnwidth}{!}{%
\begin{tabular}{lrrrrrr}
\toprule
Model & $\ell=12$ & $\ell=16$ & $\ell=20$ & $\ell=24$ & $\ell=28$ & $\ell=32$\\
\midrule
MLP/MNIST & 96.80 (6) & 97.49 (8) & 97.49 (8) & 97.49 (8) & 97.49 (8) & 97.49 (8)\\
ResNet20/CIFAR-10 & 14.84 (7) & 76.40 (8) & 87.97 (11) & 88.23 (14) & 88.23 (14) & 88.23 (14)\\
LSTM/synthetic, ideal gates & 95.60 (7) & 95.60 (7) & 95.60 (7) & 95.60 (7) & 95.60 (7) & 95.60 (7)\\
\bottomrule
\end{tabular}
}
\end{center}

The historical ImageNette E0-U-S results are shown in
Table~\ref{tab:e0-uniform}. Each model contributes all
132 calibration candidates and six validation rows. The table reports the
training-selected $s$, validation Top-1, the difference from the FP32
reference, and per-example agreement with FP32.

Under E0-U-S arithmetic, the selected ResNet-50 16-bit format reaches
$65.554\%$ (a $-9.860$
percentage-point difference), while 20 bits reaches $74.803\%$. At 24 bits,
Top-1 is $75.439\%$ versus the FP32 $75.414\%$, but agreement is $97.834\%$;
28 and 32 bits use the same $s=20$ and are exactly equal, with $99.949\%$
agreement. ResNet-101 shows $62.726\%$ at 16 bits and $78.599\%$ at 20 bits.
Its 24, 28, and 32-bit selected configurations use $s=15$, have no storage
clipping, and produce identical predictions with $78.344\%$ Top-1 and
$97.656\%$ agreement. ResNet-152 reaches $59.822\%$ at 16 bits and
$77.401\%$ at 20 bits; its 24, 28, and 32-bit rows are identical at
$77.605\%$ Top-1 and $97.707\%$ agreement.

In this saturating control, the temporary register has a fixed logical width
of 128 bits. CUDA evaluates
the linear operators with integer limb matrix multiplication and exact
128-bit reconstruction. Absolute-value bounds cover every original
accumulation prefix and reconstruction prefix; all 396 calibration and 18
validation rows complete without temporary overflow. The FP32 reference
disables TF32 and attains $75.414\%$, $78.089\%$, and $77.478\%$ for the
three models, respectively.

The saturating-control results separate two precision effects. The 16-bit rows have substantial
parameter or activation clipping, while several 20-bit rows still have a small
number of clipped writes. The selected 24-bit rows have no clipping and are
within $0.255$ percentage points of their FP32 references, but agreement remains
below 100\%. These results motivate wider fixed-point formats for preserving
accuracy in large CNNs. The required width depends on the model, scale, range,
and fidelity metric. Scale selection maximizes calibration Top-1, so the
reported rows do not optimize prediction
agreement: for ResNet-101/152, $s=15$ wins Top-1 although higher-$s$
calibration candidates attain greater FP32 agreement. Increasing $\ell$ alone
leaves the selected $s=15$ predictions unchanged. The 24-bit accuracy result
is specific to the wide-accumulator saturating control; it does not establish
the accuracy of strict $\mathbb Z_{2^\ell}$ inference, for which no complete
validation results are reported here. The logical 128-bit temporary is a
cleartext arithmetic choice, not an implemented 128-bit two-party accumulator.
Matching its accuracy in a secure backend requires the corresponding expansion,
truncation, and saturation semantics.

\begin{table}[t]
\centering
\caption{Historical cleartext E0-U-S saturating control: uniform $Q(\ell,s)$ storage for parameters, inputs, intermediate tensors, and logits, with floor-and-saturate encoding, logical 128-bit integer accumulation, and one floor rescaling followed by saturating writeback. These values do not use strict $\mathbb Z_{2^\ell}$ arithmetic. For each $\ell$, $s$ maximizes Top-1 on 512 class-balanced training images, with ties choosing smaller $s$; evaluation uses all 3,925 ImageNette validation images. $\Delta$ is Top-1 minus FP32 in percentage points, and agreement is the percentage of predictions matching FP32. The FP32 reference disables CUDA matmul and cuDNN TF32.}
\label{tab:e0-uniform}
\begingroup
\footnotesize
\setlength{\tabcolsep}{3pt}
\begin{tabular}{lrrrrr}
\toprule
Model & $\ell$ & $s$ & \shortstack{Top-1\\(\%)} & \shortstack{$\Delta$\\(pp)} & \shortstack{Agreement\\(\%)} \\
\midrule
ResNet-50 & 12 & 0 & 0.000 & -75.414 & 0.000 \\
 & 16 & 12 & 65.554 & -9.860 & 75.924 \\
 & 20 & 13 & 74.803 & -0.611 & 91.847 \\
 & 24 & 15 & 75.439 & +0.025 & 97.834 \\
 & 28 & 20 & 75.389 & -0.025 & 99.949 \\
 & 32 & 20 & 75.389 & -0.025 & 99.949 \\
\addlinespace
ResNet-101 & 12 & 0 & 0.000 & -78.089 & 0.000 \\
 & 16 & 12 & 62.726 & -15.363 & 72.331 \\
 & 20 & 14 & 78.599 & +0.510 & 95.873 \\
 & 24 & 15 & 78.344 & +0.255 & 97.656 \\
 & 28 & 15 & 78.344 & +0.255 & 97.656 \\
 & 32 & 15 & 78.344 & +0.255 & 97.656 \\
\addlinespace
ResNet-152 & 12 & 0 & 0.000 & -77.478 & 0.000 \\
 & 16 & 12 & 59.822 & -17.656 & 70.140 \\
 & 20 & 14 & 77.401 & -0.076 & 95.414 \\
 & 24 & 15 & 77.605 & +0.127 & 97.707 \\
 & 28 & 15 & 77.605 & +0.127 & 97.707 \\
 & 32 & 15 & 77.605 & +0.127 & 97.707 \\
\bottomrule
\end{tabular}
\endgroup
\end{table}


The strict E0-Q-fixed control retains its original results. Its selected
20-bit configurations all use $s=15$ and attain $71.847\%$, $69.936\%$, and
$62.548\%$ Top-1; the 24-bit configurations use $s=17$ and attain
$79.694\%$, $81.376\%$, and $79.720\%$. Per-product rescaling and
saturating accumulation yield different accuracy curves. In particular,
saturating the $7\times7$ pooling sum before division limits a nonnegative
Q$(24,17)$ pooled feature to approximately $64/49$. The historical FP32
references are $75.414\%$, $78.115\%$, and $77.478\%$; the historical
ResNet-101 baseline differs from E0-U-S by one example, and the earlier run did
not record its TF32 settings. The control measures storage quantization together with these
accumulation effects.

\paragraph*{Narrow ReLU interfaces with a wider backbone}
A separate ImageNette experiment retains Q$(24,15)$ storage outside ReLU and
wide accumulation, but converts each ReLU input and output to a signed 16-bit
format. Conversion uses signed floor rescaling and saturation, followed by an
exact shift back to the backbone scale. The pretrained weights, original
1,000-class output head, and 3,925-image validation set are unchanged.
Fractional widths are selected on 512 training images; Q$(16,12)$ is an
additional prespecified comparison.

At Q$(16,12)$, ResNet-50 reaches $75.2866\%$ Top-1 against the FP32
$75.4140\%$, and ResNet-152 reaches $77.3248\%$ against $77.4777\%$.
The corresponding FP32 prediction agreements are $97.5796\%$ and
$97.2484\%$. These losses are smaller than those of uniform 16-bit storage.
The training-selected ReLU formats are Q$(16,14)$ for ResNet-50 and
Q$(16,11)$ for ResNet-152, with Top-1 of $76.2548\%$ and $77.2484\%$.
For ResNet-50, the higher Top-1 coexists with only $90.0637\%$ FP32 agreement,
showing why the two metrics must be reported separately. These cleartext results
show that narrow activation interfaces can preserve accuracy within a wider
numerical program. Their secure cost includes rescaling and saturation as well
as the activation itself. The comparison domains of these conversions depend
on the implemented rounding and range semantics; the following trace illustrates
wide-domain comparisons in mixed-precision rescaling.

\paragraph*{Wide comparison domains in mixed-precision rescaling}
We analyze an existing single-step Google-30/FastGRNN execution with batch size
128 over ten fresh queries. Six SCI-compatible signed rescaling calls, arising
from matrix multiplication and elementwise or scalar products, retain their
24- or 32-bit output rings and round division by a power of two toward zero.
Per query, these calls invoke $6{,}528$ comparisons over 23-bit domains and
$51{,}200$ comparisons over 31-bit domains, counting each joint instance once.
Thus the mixed-width program retains substantial wide-domain nonlinear work
alongside its narrow tensor and activation interfaces. The domains follow from
these signed rescaling operations: truncation with ring reduction can instead
use smaller domains, such as a 4-bit comparison for a 36-to-32-bit reduction
with a right shift of four. The relevant FSS width is that of the actual
comparison, rather than the storage width or source ring alone.

\paragraph*{Historical Google-30 replay}
We reproduced the SiRNN toolchain using Speech Commands v0.01, the official
EdgeML FastGRNN implementation, SeeDot quantization, and SCI-OT. The weights
were trained locally with seed 42, hidden width 100, low-rank dimensions 16/35,
and 300 training epochs; they are not the paper authors' original weights.
The generated program processes 99 steps of 32 features and predicts 30 classes.
It uses mixed 8/16/32-bit values rather than a uniform 32-bit numerical format.
Its sigmoid interface uses 16-bit input/output with fractional widths 7 and 14;
tanh uses 16-bit input/output with fractional width 7.
The server privately holds matrix weights and biases, while the learned
scalars $\zeta$ and $\nu$ are public constants in the compiled program.

On the 6,835-example test split, the historical FP32 model obtains
$6{,}177/6{,}835=90.3731\%$ and the SeeDot reference obtains
$6{,}134/6{,}835=89.7440\%$. These are historical reference accuracies:
checkpoint selection used test accuracy, and the input-bound computation used
both training and test data. They therefore do not constitute an independent
final quality evaluation.

SCI was run on one fixed example per class, for 30 complete sequences.
All 30 predictions match the SeeDot reference; 27 match the labels.
This establishes final-prediction consistency on that subset, not full-test-set
secure accuracy or per-operator integer equivalence. In particular, the
historical SeeDot reference uses floating exponential/tanh evaluation before
integer encoding, whereas SCI uses its secure function approximations.
The SCI replay also extends 8-bit logits to 16 bits before ArgMax to preserve
the signed comparison range. Its cross-sample latency summaries are replay
statistics, not repeated system-performance estimates.

\subsection{Efficiency of the Input Decomposition}\label{sec:eval-split}

The measured segmented comparison benchmark uses two parties on localhost,
32-bit inputs, and ten measured runs for each segment count.
Table~\ref{tab:seg-measured} reports run means. Offline work includes arithmetic
threshold bit decomposition, batched cGGM/DPF generation, correction-word
exchange, and Boolean triples for the cross-segment scan. The online phase
accepts a public query and includes local tree evaluation and the scan's masked
openings. Private-input masking and final output reconstruction are outside
this timing boundary. Consequently these rows measure a comparison primitive,
not private network inference.

\begin{table}[t]
\centering
\caption{Measured 32-bit two-party localhost comparison primitive. Times are
means of ten runs. Online bytes are the historical sent-plus-received counter
of one party (link traffic), and rounds count dependent scan openings.
The query is public; private-input conversion and final reconstruction are
excluded.}
\label{tab:seg-measured}
\small
\setlength{\tabcolsep}{4pt}
\begin{tabular}{rrrrrr}
\toprule
$\nseg$ & \shortstack{Offline\\ms} & \shortstack{Online\\ms} &
\shortstack{Total\\ms} & \shortstack{Online\\bytes} & \shortstack{Scan\\rounds}\\
\midrule
2  & 39.8306 & 0.2208 & 40.0514 & 8   & 1\\
3  & 37.2425 & 0.1923 & 37.4348 & 16  & 2\\
4  & 36.3243 & 0.2240 & 36.5483 & 24  & 2\\
8  & 40.5740 & 0.4577 & 41.0317 & 56  & 3\\
16 & 40.0162 & 1.5401 & 41.5563 & 120 & 4\\
\bottomrule
\end{tabular}
\end{table}

All measured rows pass comparison reconstruction. Among these configurations,
$\nseg=4$ has the lowest mean total time, $36.5483$\,ms, with $277{,}920$
offline bytes under the same counter convention. This is a local measurement,
not a LAN/WAN optimum. No completed depth-32 monolithic row is available, so
we do not assign it a latency or derive a measured speedup from it. The online
scan depth is $\lceil\log_2\nseg\rceil$, consistent with the balanced reduction.
The measurements also charge arithmetic threshold decomposition; the tree-only
bound in Theorem~\ref{thm:cost} starts from Boolean-shared threshold bits.

\subsection{Planned Comparison on Complete Integer Programs}\label{sec:eval-compose}

The next comparison will start with the native Google-30 mixed-width program.
The first target backend, \emph{SCI + ours}, will preserve SCI's linear layers,
expansion, truncation, activation approximations, and output reconstruction.
Only matching comparison/MSB operations will be replaced by segmented FSS.
Acceptance requires exact agreement of reconstructed integer operator outputs
with SCI and identical final predictions on the complete fixed test set;
random secret shares themselves need not match. This same-program test differs
from the historical SeeDot-to-SCI final-prediction replay. The separate
SIGMA-based activation in Section~\ref{sec:infer} will not be substituted into
this comparison because it changes the numerical algorithm.

The system comparison will allow each backend to use its native width
optimizations subject to a prespecified Top-1 loss of at most 1.0 percentage
point relative to the same FP32 checkpoint. Prediction agreement will be
reported alongside Top-1. Checkpoint and quantization choices will use training
or calibration data before the final evaluation. A second comparison will hold
the complete integer program fixed, including scales, accumulation widths,
rounding, truncation, and activation approximations, to isolate backend costs.
SCI/SiRNN OT, SCI + ours, and MP-SPDZ~\cite{mpspdz} \texttt{semi2k} with real OT preprocessing
will form the two-party, dealerless, semi-honest comparison once the matching
executors are available. ABY-64 will provide uniform-64-bit context; its separate
placement does not imply a trusted dealer. Dealer-assisted FSS will be reported
under its additional preprocessing trust assumption.

Width controls will compare the native program with uniform 16-bit and 32-bit
storage/I/O, while allowing necessary wider products and accumulators and
per-tensor fractional scales. A strict same-ring variant will diagnose modular
wraparound separately. Uniform 64-bit storage is a conditional extension.
The experiment will record actual comparison/FSS domain widths, invocation
counts, and conversion costs, and compare direct wide-domain computation with
rescaling, narrow-domain computation, and any necessary restoration. These
controls test width requirements for the evaluated models and quality criterion,
not a universal lower bound.

Performance runs will use identical hardware, thread counts, batches, and
input order, with LAN/WAN conditions and repetitions fixed before measurement.
Results will include offline, online, and total latency; communication as both
parties' sent-byte sum; dependent rounds; and each party's peak memory.
Initialization, private-input conversion, key/triple generation, expansion,
truncation, and reconstruction will have explicit phase boundaries. Any
initialization amortization will state its query count, while one-time keys,
masks, and triples will be freshly consumed for each query. A width/segment-count
sweep will include the domains used by the program and 16/20/24/32-bit coverage;
timeouts and memory limits will remain censored observations. For CNNs, the
end-to-end comparison will use ResNet-50 and ResNet-152 with a uniform
Q$(32,12)$ configuration, accounting for conversions alongside nonlinear
evaluation. Secure end-to-end performance for these segmented-backend networks
has not yet been measured.

% =====================================================================
\section{Conclusion}\label{sec:conc}

Segmenting the function domain reduces the exponential tree-expansion cost
of dealerless FSS while retaining a balanced secure scan for comparison.
Two-party localhost measurements establish a working 32-bit comparison
primitive. The numerical studies motivate wide fixed-point arithmetic to
preserve accuracy in large networks, while the mixed-width FastGRNN trace
identifies 23- and 31-bit comparison domains in signed rescaling. Together,
these results motivate efficient dealerless generation for wide nonlinear
domains even when tensor storage and activation interfaces use mixed precision.
The historical Google-30 replay supplies subset-level SCI consistency evidence.
The next step is to measure whether the segmented comparison improves a
complete mixed-width program under fixed integer
semantics and a shared accuracy threshold; the present results do not establish
an end-to-end advantage over OT inference.

% =====================================================================
\appendices

\section{Full Generation and Evaluation Procedures}\label{app:full}

The per-segment generation $\Pi_{\mathsf{Seg}}$ and evaluation are both given in
full in the main body (Algorithm~\ref{alg:gen}), building on $\Pi_{\mathsf{Split}}$
(Algorithm~\ref{alg:split}) and $\PREP$ (Algorithm~\ref{alg:prep}). The \textbf{[Offline]}
block constructs the boolean DPF key via Half-Tree correlated GGM with parallel COT
correlations. The \textbf{[Online]} block shows the Grotto prefix-parity walk that
recovers the boolean comparison share from the opened $\xhat\seg{j}$ with $\bigO(\wseg_j)$
node touches and one $\BitBA$ mixed product. Here we detail how the cross-segment
scan completes the overall comparison from the $\nseg$ per-segment boolean shares.

% =====================================================================


\section{FSS-based Nonlinear Building Blocks}\label{appendix:Secure 3PC Protocols}

\subsection{GELU and SiLU}
To reduce overhead, we design a secure GELU protocol ($\Pi_{\text{GELU}}$,
Algorithm~\ref{alg:gelu}) inspired by SIGMA~\cite{sigma}. Because the encoded
difference $\delta(x)=\text{ReLU}(x)-\text{GELU}(x)$ is rounded to zero
outside $[-4,4]$, we tabulate only $[0,4]$. Discretizing this interval at
$2^{-6}$ precision gives 256 entries,
$\vec{T}_{\text{GELU}}[i]=\lfloor\delta(i/2^6)2^f\rfloor$ for
$i\in[0,255]$. SiLU uses the same LUT-oriented pipeline with a different
public table~\cite{sigma}.

\begin{algorithm}[ht]
  \caption{GELU ($\Pi_{\text{GELU}}$)}
  \label{alg:gelu}
  \textbf{Input:} $P_b$ holds $\ashare{x}_b$ and the public table $\vec{T}_{\text{GELU}}$. \\
  \textbf{Output:} $P_b$ holds $\ashare{z}_b$ where $z=\text{GELU}(x)$.
  \begin{algorithmic}[1]
    \State $\ashare y \leftarrow \Trunc(\ashare x,f-6)$
    \State $\ashare{p_x}\leftarrow\Pi_{\mathsf{Select}}(\ashare y,\ashare x)$ and $\ashare{p_y}\leftarrow\Pi_{\mathsf{Select}}(\ashare y,\ashare y)$, reusing the sign bit of $y$
    \State $\ashare a\leftarrow2\ashare{p_y}-\ashare y$, $\ashare\delta\leftarrow256-\ashare a$
    \State $\ashare i\leftarrow\Pi_{\mathsf{Select}}(\ashare\delta,\ashare a)$
    \State \textbf{return} $\ashare z=\ashare{p_x}-\Pi_{\mathsf{LUT}}(\ashare i,\vec{T}_{\text{GELU}})$
  \end{algorithmic}
\end{algorithm}

\textbf{Complexity.} The two Select calls in Line~2 reuse one DReLU bit but
require two $\ell$-bit $\BitBA$ calls. The later Select is independent, and
the table index has $q_g=8$ bits. Consequently,
\begin{equation}
\begin{aligned}
C^{\mathrm{off}}_{\mathsf{GELU}}
 &=C^{\mathrm{off}}_{\mathsf{Trunc}}(f-6)
   +C^{\mathrm{off}}_{\DReLU}\\
 &\quad+8\ell(\lambda+\ell+1)+C^{\mathrm{off}}_{\ReLU}\\
 &\quad+C^{\mathrm{off}}_{\mathsf{LUT}}(8,\ell),\\
R^{\mathrm{off}}_{\mathsf{GELU}}
 &=\max\{R^{\mathrm{off}}_{\mathsf{Trunc}}(f-6),
          R^{\mathrm{off}}_{\ReLU},\\
 &\qquad R^{\mathrm{off}}_{\mathsf{LUT}}(8,\ell),3\},\\
C^{\mathrm{on}}_{\mathsf{GELU}}
 &=C^{\mathrm{on}}_{\mathsf{Trunc}}(f-6)
   +C^{\mathrm{on}}_{\DReLU}+2(\ell+1)
   +C^{\mathrm{on}}_{\ReLU}+8,\\
R^{\mathrm{on}}_{\mathsf{GELU}}
 &=R^{\mathrm{on}}_{\mathsf{Trunc}}(f-6)
   +2R^{\mathrm{on}}_{\ReLU}+1.
\end{aligned}
\label{eq:gelu-cost}
\end{equation}
The online stages are sequential because each produces the next index. SiLU
has the same cost when instantiated with the stated pipeline.

\subsection{Inverse and Reciprocal Square Root}
Following Mosformer~\cite{Mosformer}, we compute division
($\Pi_{\mathsf{Inv}}$, Algorithm~\ref{alg:inverse}) via compressed intervals.
A positive input $x$ is normalized into $[1/b,1]$ by finding $k$ such that
$b^k\leq x<b^{k+1}$. We apply $\Pi_{\DReLU}$ against the $n$ public
thresholds $b^i$ in one batch, where
$n=\lfloor\log_b(2^{2f})-1\rfloor$, convert the resulting bits to arithmetic
shares, and set $k=\sum_i d[i]$. A first LUT returns the factor
$b^{-(k+1)}$ and a second LUT evaluates the reciprocal on the normalized
domain. The latter has $(1-1/b)2^f$ entries before power-of-two padding.

Our reciprocal square-root protocol ($\Pi_{\mathsf{Rsqrt}}$,
Algorithm~\ref{alg:rsqrt}) reuses this interval-identification pipeline. It
adds a parallel scaling-factor LUT and substitutes the reciprocal-square-root
table in the normalized domain.

\begin{algorithm}[ht]
  \caption{Inverse ($\Pi_{\text{Inv}}$)}
  \label{alg:inverse}
  \textbf{Input:} $P_b$ holds $\ashare{x}_b$, scaling base $b$, precision $f$, and public tables $\vec{T}_{\text{exp}}$ and $\vec{T}_{\text{inv}}$. \\
  \textbf{Output:} $P_b$ holds $\ashare{z}_b$ where $z=1/x$.
  \begin{algorithmic}[1]
    \For{$i = 1$ to $n$ in parallel}
      \State $\bshare{d[i]}\leftarrow\Pi_{\DReLU}(\ashare x-b^i)$
    \EndFor
    \State $\ashare{k}\leftarrow\sum_{i=1}^{n}\mathsf{B2A}(\bshare{d[i]})$
    \State $\ashare t\leftarrow\Pi_{\mathsf{LUT}}(\ashare{k+1},\vec{T}_{\text{exp}})$
    \State $\ashare q\leftarrow\MulTr(\ashare x,\ashare t;f)$
    \State $\ashare w\leftarrow\Pi_{\mathsf{LUT}}(\ashare{q-1/b},\vec{T}_{\text{inv}})$
    \State \textbf{return} $\ashare z\leftarrow\MulTr(\ashare t,\ashare w;f)$
  \end{algorithmic}
\end{algorithm}

\begin{algorithm}[ht]
  \caption{Reciprocal Square Root ($\Pi_{\text{Rsqrt}}$)}
  \label{alg:rsqrt}
  \textbf{Input:} As in $\Pi_{\mathsf{Inv}}$, with public tables $\vec{T}_{\text{exsqrt}}$ and $\vec{T}_{\text{rsqrt}}$. \\
  \textbf{Output:} $P_b$ holds $\ashare z_b$ where $z=1/\sqrt{x}$.
  \begin{algorithmic}[1]
    \State Perform the batched interval comparisons and B2A conversion of Algorithm~\ref{alg:inverse} to obtain $\ashare k$.
    \State $\ashare t\leftarrow\Pi_{\mathsf{LUT}}(\ashare{k+1},\vec{T}_{\text{exp}})$ and $\ashare s\leftarrow\Pi_{\mathsf{LUT}}(\ashare{k+1},\vec{T}_{\text{exsqrt}})$ in parallel
    \State $\ashare q\leftarrow\MulTr(\ashare x,\ashare t;f)$
    \State $\ashare w\leftarrow\Pi_{\mathsf{LUT}}(\ashare{q-1/b},\vec{T}_{\text{rsqrt}})$
    \State \textbf{return} $\ashare z\leftarrow\MulTr(\ashare s,\ashare w;f)$
  \end{algorithmic}
\end{algorithm}

\textbf{Complexity.} Let $q_k=\lceil\log_2(n+1)\rceil$ and let
$q_i=\lceil\log_2((1-1/b)2^f)\rceil$ after padding the normalized table. The
$n$ comparisons and conversions are parallel; the two multiplications each
include exact truncation. Inverse therefore costs
\begin{equation}
\begin{aligned}
C^{\mathrm{off}}_{\mathsf{Inv}}
 &=nC^{\mathrm{off}}_{\DReLU}
   +C^{\mathrm{off}}_{\mathsf{LUT}}(q_k,\ell)
   +C^{\mathrm{off}}_{\mathsf{LUT}}(q_i,\ell)
   +2C^{\mathrm{off}}_{\mathsf{MT}}(f),\\
R^{\mathrm{off}}_{\mathsf{Inv}}
 &=\max\{R^{\mathrm{off}}_{\DReLU},
          R^{\mathrm{off}}_{\mathsf{LUT}}(q_k,\ell),\\
 &\qquad R^{\mathrm{off}}_{\mathsf{LUT}}(q_i,\ell),
          R^{\mathrm{off}}_{\mathsf{MT}}(f)\},\\
C^{\mathrm{on}}_{\mathsf{Inv}}
 &=nC^{\mathrm{on}}_{\DReLU}+n(\lambda+\ell)+q_k+q_i
   +2C^{\mathrm{on}}_{\mathsf{MT}}(f),\\
R^{\mathrm{on}}_{\mathsf{Inv}}
 &=R^{\mathrm{on}}_{\DReLU}+2R^{\mathrm{on}}_{\mathsf{MT}}(f)+3.
\end{aligned}
\label{eq:inverse-cost}
\end{equation}
The three added rounds are the batched B2A conversion and the two dependent
LUT openings. Rsqrt evaluates two $q_k$-bit LUTs in parallel, instead of one,
and otherwise has the same dependency graph:
\begin{equation}
\begin{aligned}
C^{\mathrm{off}}_{\mathsf{Rsqrt}}
 &=nC^{\mathrm{off}}_{\DReLU}
   +2C^{\mathrm{off}}_{\mathsf{LUT}}(q_k,\ell)
   +C^{\mathrm{off}}_{\mathsf{LUT}}(q_i,\ell)
   +2C^{\mathrm{off}}_{\mathsf{MT}}(f),\\
R^{\mathrm{off}}_{\mathsf{Rsqrt}}
 &=R^{\mathrm{off}}_{\mathsf{Inv}},\\
C^{\mathrm{on}}_{\mathsf{Rsqrt}}
 &=nC^{\mathrm{on}}_{\DReLU}+n(\lambda+\ell)+2q_k+q_i
   +2C^{\mathrm{on}}_{\mathsf{MT}}(f),\\
R^{\mathrm{on}}_{\mathsf{Rsqrt}}
 &=R^{\mathrm{on}}_{\mathsf{Inv}}.
\end{aligned}
\label{eq:rsqrt-cost}
\end{equation}

\paragraph*{LayerNorm}
For a vector of width $N$, we represent the public reciprocal $1/N$ and the
affine scale at precision $f$. Computing the mean and variance average then
uses two scalar exact truncations. The variance stage also evaluates $N$
parallel square-and-truncate operations, followed by one Rsqrt; normalization
uses $N$ parallel multiply-and-truncate operations, and applying the public
affine scale uses $N$ parallel truncations. The public bias is local. Hence
\begin{equation}
\begin{aligned}
C^{\mathrm{off}}_{\mathsf{LayerNorm}}(N)
 &=2N C^{\mathrm{off}}_{\mathsf{MT}}(f)
   +(N+2)C^{\mathrm{off}}_{\mathsf{Trunc}}(f)
   +C^{\mathrm{off}}_{\mathsf{Rsqrt}},\\
R^{\mathrm{off}}_{\mathsf{LayerNorm}}(N)
 &=\max\{R^{\mathrm{off}}_{\mathsf{MT}}(f),
          R^{\mathrm{off}}_{\mathsf{Trunc}}(f),R^{\mathrm{off}}_{\mathsf{Rsqrt}}\},\\
C^{\mathrm{on}}_{\mathsf{LayerNorm}}(N)
 &=2N C^{\mathrm{on}}_{\mathsf{MT}}(f)
   +(N+2)C^{\mathrm{on}}_{\mathsf{Trunc}}(f)
   +C^{\mathrm{on}}_{\mathsf{Rsqrt}},\\
R^{\mathrm{on}}_{\mathsf{LayerNorm}}(N)
 &=2R^{\mathrm{on}}_{\mathsf{MT}}(f)
   +3R^{\mathrm{on}}_{\mathsf{Trunc}}(f)
   +R^{\mathrm{on}}_{\mathsf{Rsqrt}}.
\end{aligned}
\label{eq:layernorm-cost}
\end{equation}
The factor $N$ affects communication but not the depth of a batched vector
stage. If $1/N$ or the affine scale is exactly representable without
rescaling, the corresponding truncation term is omitted.

\subsection{Softmax}
$\Pi_{\mathsf{Softmax}}$ (Algorithm~\ref{alg:softmax}) composes MaxPool,
Exponentiation, and Inverse. A balanced MaxPool tree over $N$ inputs has
$N-1$ internal nodes and depth $\lceil\log_2N\rceil$; exponentiations and the
final products are each batched across all coordinates.

\begin{algorithm}[ht]
  \caption{Softmax ($\Pi_{\text{Softmax}}$)}
  \label{alg:softmax}
  \textbf{Input:} $P_b$ holds $\ashare{\vec{x}}_b$ of length $N$. \\
  \textbf{Output:} $P_b$ holds $\ashare{\vec{z}}_b$ where $\vec{z}=\text{Softmax}(\vec{x})$.
  \begin{algorithmic}[1]
    \State $\ashare{x_{\max}}\leftarrow\Pi_{\mathsf{MaxPool}}(\ashare{\vec{x}})$ using a balanced tree.
    \For{$i = 0$ to $N-1$ in parallel}
      \State $\ashare{\delta[i]}\leftarrow\ashare{\vec{x}[i]}-\ashare{x_{\max}}$
      \State $\ashare{\tilde e[i]}\leftarrow\Pi_{\mathsf{Exp}}(\ashare{\delta[i]})$
    \EndFor
    \State $\ashare s\leftarrow\sum_{i=0}^{N-1}\ashare{\tilde e[i]}$, $\ashare r\leftarrow\Pi_{\mathsf{Inv}}(\ashare s)$
    \For{$i = 0$ to $N-1$ in parallel}
      \State $\ashare{\vec z[i]}\leftarrow\MulTr(\ashare{\tilde e[i]},\ashare r;f)$
    \EndFor
    \State \textbf{return} $\ashare{\vec z}$
  \end{algorithmic}
\end{algorithm}

\textbf{Complexity.} Let $q_e$ be the exponentiation-table index width and
use the Inverse parameters of~\eqref{eq:inverse-cost}. The MaxPool precedes
the exponentiations, which precede Inverse and the final products. Therefore
\begin{equation}
\begin{aligned}
C^{\mathrm{off}}_{\mathsf{Softmax}}(N)
 &=C^{\mathrm{off}}_{\mathsf{MaxPool}}(N)
   +NC^{\mathrm{off}}_{\mathsf{Exp}}(q_e)\\
 &\quad+C^{\mathrm{off}}_{\mathsf{Inv}}
   +NC^{\mathrm{off}}_{\mathsf{MT}}(f),\\
R^{\mathrm{off}}_{\mathsf{Softmax}}(N)
 &=\max\{R^{\mathrm{off}}_{\mathsf{MaxPool}},
          R^{\mathrm{off}}_{\mathsf{Exp}}(q_e),\\
 &\qquad R^{\mathrm{off}}_{\mathsf{Inv}},
          R^{\mathrm{off}}_{\mathsf{MT}}(f)\},\\
C^{\mathrm{on}}_{\mathsf{Softmax}}(N)
 &=C^{\mathrm{on}}_{\mathsf{MaxPool}}(N)
   +NC^{\mathrm{on}}_{\mathsf{Exp}}(q_e)\\
 &\quad+C^{\mathrm{on}}_{\mathsf{Inv}}
   +NC^{\mathrm{on}}_{\mathsf{MT}}(f),\\
R^{\mathrm{on}}_{\mathsf{Softmax}}(N)
 &=R^{\mathrm{on}}_{\mathsf{MaxPool}}(N)
   +R^{\mathrm{on}}_{\mathsf{Exp}}(q_e)
   +R^{\mathrm{on}}_{\mathsf{Inv}}+R^{\mathrm{on}}_{\mathsf{MT}}(f).
\end{aligned}
\label{eq:softmax-cost}
\end{equation}
Thus $N$ multiplies communication in the two elementwise stages, while their
round depth is unchanged by batching. The only logarithmic vector-length term
is the MaxPool depth.


% =====================================================================
\bibliographystyle{IEEEtran}
\bibliography{references}

\end{document}
